% =====================================================================
%  第 8 章 中心与 opers 上函数代数的等同
%  §8.1（含 8.1.1–8.1.5）、§8.2（含 8.2.1–8.2.6）、§8.3（含 8.3.1–8.3.3）
%  对应 PDF p.232–p.260；页脚印刷页码 p.221–p.249
%  （任务书标注“印刷 p.218–246”系按 PDF−14 推算，实际页脚为 PDF−11，
%    本文件批注中的 p.N 一律取页脚上的印刷页码。）
%  编号公式：(8.1-1)–(8.1-11)、(8.2-1)–(8.2-19)、(8.3-1)–(8.3-4)
% =====================================================================

\chapter{中心与 opers 上函数代数的等同}

本章利用前面把顶点代数 $V_{\kappa_c}(\mathfrak g)$ 的中心写成屏蔽算子核的形式这一描述，完成 §7.1 勾画的计划。论证遵循 [Frenkel 2005b]。我们先证明，该中心典范同构于 Langlands 对偶李代数 ${}^L\mathfrak g$ 对应的所谓经典 W-代数。接着证明，经典 W-代数又同构于形式圆盘 $D$ 上 ${}^LG$-opers 空间的函数代数。$D$ 上 ${}^LG$-opers 空间的函数代数 $\operatorname{Fun}\operatorname{Op}_{{}^LG}(D)$ 带有顶点 Poisson 结构。我们证明，$V_{\kappa_c}(\mathfrak g)$ 的中心（带由水平形变给出的典范顶点 Poisson 结构）与 $\operatorname{Fun}\operatorname{Op}_{{}^LG}(D)$ 之间的这个同构保持顶点 Poisson 结构。借助该同构，我们证明临界水平下 $\hat{\mathfrak g}$ 的完备泛包络代数的中心，作为 Poisson 代数同构于去心圆盘 $D^\times$ 上 opers 空间的函数代数 $\operatorname{Fun}\operatorname{Op}_{{}^LG}(D^\times)$。后一同构由 V. Drinfeld 猜想。

下面是本章内容的更详细说明。在 §8.1 中，我们把顶点代数 $V_{\kappa_c}(\mathfrak g)$ 的中心与若干算子的核之交等同起来，并进而把该交与 Langlands 对偶李代数 ${}^L\mathfrak g$ 对应的经典 W-代数等同。我们还证明相应的顶点 Poisson 代数结构一致。在 §8.2.1 中我们定义 Miura opers，并解释它们与某个 $H$-丛上的联络的关系。随后在 §8.2 中证明：这个经典 W-代数（相应地，交换顶点代数 $\pi_0$）不是别的，正是圆盘上 ${}^LG$-opers 空间（相应地，圆盘上 Miura ${}^LG$-opers 空间）上的函数代数。此外，经典 W-代数到 $\pi_0$ 的嵌入，与两个函数代数之间的 Miura 映射一致。最后，在 §8.3 中我们考虑临界水平下 $\hat{\mathfrak g}$ 的完备泛包络代数的中心，把它与去心圆盘上 ${}^LG$-opers 空间的函数代数等同，并证明该等同满足种种相容性。特别地，我们证明：把中心元在 Wakimoto 模上取值所得的 Harish--Chandra 同态的仿射类比，正是从去心圆盘上 Miura ${}^LG$-opers 到 ${}^LG$-opers 的 Miura 变换。

\begin{keypoints}
    \item \textbf{核心思想：}$V_{\kappa_c}(\mathfrak g)$ 的中心 $\mathfrak z(\hat{\mathfrak g})$ 有两个面貌：一是 $\pi_0$ 上屏蔽算子 $\overline V_i[1]$ 的核之交（$\hat{\mathfrak g}$ 一侧），二是 ${}^L\mathfrak g$ 的经典 W-代数 $\mathbf W({}^L\mathfrak g)$（对偶一侧）。
    \item \textbf{主线：}全书计划在第 8 章收口，链条为 $\mathfrak z(\hat{\mathfrak g})\simeq\mathbf W({}^L\mathfrak g)\simeq\operatorname{Fun}\operatorname{Op}_{{}^LG}(D)$，每一步都保持顶点 Poisson 结构与 $\operatorname{Der}\mathcal O$、$\operatorname{Aut}\mathcal O$-作用。
    \item \textbf{几何工具：}Miura oper 是 oper 加一个被联络保持的第二 Borel 约化；其空间同构于 $H$-丛 $\Omega^{\check\rho}$ 上的联络空间，而 Miura 变换就是遗忘映射。
    \item \textbf{关键结论：}定理 8.2.5 给出圆盘 $D$ 上的等同，定理 8.3.1 给出去心圆盘 $D^\times$ 上作为 Poisson 代数的等同（Drinfeld 猜想）。
    \item \textbf{记号提醒：}$\pi_0(\mathfrak g)$、$\pi^\vee_0(\mathfrak g)$、$\mathbf W(\mathfrak g)$ 都是相应单参数族在 $\epsilon\to0$ 处的特殊化，即交换顶点代数；其顶点 Poisson 结构正是由该极限过程遗留下来的。
\end{keypoints}

% [Tutor Note: p.221 章首给出全章路线，请对照 §7.1 的六步计划读：第 5 步（§8.1）把中心等同于屏蔽算子核之交，第 6 步（§8.2）把该交等同于 opers 上函数代数；§8.3 再把结论从 $\mathfrak z(\hat{\mathfrak g})$ 提升到完备泛包络代数的中心 $Z(\hat{\mathfrak g})$，并用 $\widetilde U$ 函子把顶点 Poisson 代数变成拓扑 Poisson 代数。这条链条的终点是定理 4.3.2，因此本章本质上是在补证第 4 章引用的那个结论。]

\section{顶点代数中心的描述}

\begin{keypoints}
    \item \textbf{核心思想：}先用特征标把 $\mathfrak z(\hat{\mathfrak g})$ 与屏蔽算子核之交“夹”到相等，再把这个交认成 Langlands 对偶李代数的经典 W-代数。
    \item \textbf{数学工具：}PBW 滤过与分次伴随、$\tilde{\mathfrak b}_+$-不变量的射流描述（引理 8.1.1）、$\operatorname{Inv}\mathfrak g^*[\![t]\!]$ 的自由多项式结构（定理 3.4.2）。
    \item \textbf{关键结论：}定理 8.1.5（中心 $\simeq\mathbf W({}^L\mathfrak g)$）与定理 8.1.7（该同构保持顶点 Poisson 结构，并与 $\operatorname{Aut}\mathcal O$ 作用交换）。
    \item \textbf{记号提醒：}$\mathbf b'_{i,n}=\epsilon\frac{2}{\upsilon_0(h_i,h_i)}\mathbf b_{i,n}$ 中的 $\epsilon$ 是 $\upsilon_0$ 与 $\upsilon$ 之比；换基时出现的负号 $b_{i,n}\mapsto-\mathbf b'_{i,n}$ 来自根与余根的对换，不是笔误。
\end{keypoints}

本节用命题 7.3.6 描述 §3.3.2 定义的 $V_{\kappa_c}(\mathfrak g)$ 的中心 $\mathfrak z(\hat{\mathfrak g})$。

按命题 7.3.6，$\mathfrak z(\hat{\mathfrak g})$ 等同于 $\pi_0$ 中算子 $\overline V_i[1]$（$i=1,\dots,\ell$）的核之交里的一个子空间。现在证明 $\mathfrak z(\hat{\mathfrak g})$ 实际上就等于这个交。这就是 §7.1 所述计划中的第 5 步。

\subsection{$\mathfrak z(\hat{\mathfrak g})$ 的特征标计算}

如 §3.3.3 所述，$\hat{\mathfrak g}_{\kappa_c}$-模 $V_{\kappa_c}(\mathfrak g)$ 带有由泛包络代数 $U(\hat{\mathfrak g}_{\kappa_c})$ 上的 Poincaré–Birkhoff–Witt 滤过诱导的自然滤过，其分次伴随空间 $\operatorname{gr}V_{\kappa_c}(\mathfrak g)$ 同构于
\[
    \operatorname{Sym}\bigl(\mathfrak g((t))/\mathfrak g[\![t]\!]\bigr)\simeq\operatorname{Fun}\mathfrak g^*[\![t]\!],
\]
这里用到下述（依赖坐标的）配对
\begin{equation}
    \langle A\otimes f(t),\,B\otimes g(t)\rangle=\langle A,B\rangle\operatorname{Res}_{t=0}f(t)g(t)\,dt
    \tag{8.1-1}
\end{equation}
其中 $A\in\mathfrak g^*$，$B\in\mathfrak g$。

回想 $\mathfrak z(\hat{\mathfrak g})$ 等于 $V_{\kappa_c}(\mathfrak g)$ 中 $\mathfrak g[\![t]\!]$-不变向量构成的空间。$V_{\kappa_c}(\mathfrak g)$ 中 $\mathfrak g[\![t]\!]$-不变向量的符号，是 $\operatorname{gr}V_{\kappa_c}(\mathfrak g)$ 中的 $\mathfrak g[\![t]\!]$-不变向量，即 $\operatorname{Fun}\mathfrak g^*[\![t]\!]$ 中 $\mathfrak g[\![t]\!]$-不变向量空间的元素。于是得到单射
\begin{equation}
    \operatorname{gr}\mathfrak z(\hat{\mathfrak g})\hookrightarrow(\operatorname{Fun}\mathfrak g^*[\![t]\!])^{\mathfrak g[\![t]\!]}=\operatorname{Inv}\mathfrak g^*[\![t]\!]
    \tag{8.1-2}
\end{equation}
（见引理 3.3.1）。

按定理 3.4.2，代数 $\operatorname{Inv}\mathfrak g^*[\![t]\!]$ 是以 $\overline P_{i,n}$（$i=1,\dots,\ell$；$n<0$）为生成元的自由多项式代数（多项式 $\overline P_{i,n}$ 的定义见 §3.3.4）。\transnote{公式 (8.1-1) 的留数配对使 $\mathfrak g((t))$ 与 $\mathfrak g^*((t))$ 对偶，坐标 $\overline J^a_n$ 只拾取 $t^{-n-1}dt$ 的留数，故当 $A(t)\in\mathfrak g^*[\![t]\!]$ 时只有 $n<0$ 非零；这就是此处只取 $n<0$ 的 $\overline P_{i,n}$、而 (8.1-3) 的乘积从 $n_i\ge d_i+1$ 起算的原因。半开条件「$n<0$，或 $n=0$ 且 $J^a\in\mathfrak n_-$」把多出的一层方向记到正指标一侧，指标换算须按此进行。落点：§8.1.1，(8.1-1)–(8.1-4)，原书 p.222–223。[来源：本书 p.222–224；Frenkel–Ben-Zvi, \emph{Vertex Algebras and Algebraic Curves}, 2nd ed., AMS 2004, Ch. 18.]}

算子 $L_0=-t\partial_t$ 作用在 $\operatorname{Fun}\mathfrak g^*[\![t]\!]$ 上。它给出 $\operatorname{Fun}\mathfrak g^*[\![t]\!]$ 上的 $\mathbb Z$-分次，使 $\deg\overline J^a_n=-n$。于是 $\deg\overline P_{i,n}=d_i-n$；又按公式 (4.3-3)，$\operatorname{Inv}\mathfrak g^*[\![t]\!]$ 的特征标等于
\begin{equation}
    \operatorname{ch}\operatorname{Inv}\mathfrak g^*[\![t]\!]=\prod_{i=1}^{\ell}\prod_{n_i\ge d_i+1}(1-q^{n_i})^{-1}.
    \tag{8.1-3}
\end{equation}

% [Tutor Note: p.222 式 (8.1-1) 与 (8.1-3) 的指标约定极易混淆。(8.1-1) 的配对把 $\mathfrak g^*((t))$ 与 $\mathfrak g((t))$ 配对（留数配对），故 $\operatorname{Fun}\mathfrak g^*[\![t]\!]$ 的坐标是 $\overline J^a_n$，$n<0$，对应 $t^{-n-1}$，即 $t^{0},t^{1},\dots$；(8.1-3) 中的乘积指标是 $n_i\ge d_i+1$ 而不是 $n_i\ge1$，因为生成元 $\overline P_{i,n}$ 的次数为 $d_i-n$。若把 $d_i$ 与 $h^\vee$ 或 $\ell$ 混用，此处必然算错特征标。]

现在用命题 6.3.3 证明映射 (8.1-2) 是同构。考虑 $\mathfrak g[\![t]\!]$ 的李子代数 $\tilde{\mathfrak b}_+=(\mathfrak b_+\otimes1)\oplus(\mathfrak g\otimes t\mathbb C[[t]])$。自然满同态 $\mathbb M_{0,\kappa_c}\to V_{\kappa_c}(\mathfrak g)$ 给出相应 $\tilde{\mathfrak b}_+$-不变子空间之间的映射
\[
    \phi:(\mathbb M_{0,\kappa_c})^{\tilde{\mathfrak b}_+}\longrightarrow V_{\kappa_c}(\mathfrak g)^{\tilde{\mathfrak b}_+}.
\]
$\mathbb M_{0,\kappa_c}$ 与 $V_{\kappa_c}(\mathfrak g)$ 都带有自然滤过（由 $U_{\kappa_c}(\hat{\mathfrak g})$ 上的 PBW 滤过诱导），且二者之间的同态保持这些滤过。因此有分次伴随空间之间的相应映射
\[
    \phi_{\mathrm{cl}}:(\operatorname{gr}\mathbb M_{0,\kappa_c})^{\tilde{\mathfrak b}_+}\longrightarrow(\operatorname{gr}V_{\kappa_c}(\mathfrak g))^{\tilde{\mathfrak b}_+}.
\]

由于 $V_{\kappa_c}(\mathfrak g)$ 是 $\mathfrak g[\![t]\!]$ 的常值子代数 $\mathfrak g\otimes1$ 的有限维表示的直和，$V_{\kappa_c}(\mathfrak g)$ 中任一 $\tilde{\mathfrak b}_+$-不变向量都自动是 $\mathfrak g[\![t]\!]$-不变的。因此
\[
    V_{\kappa_c}(\mathfrak g)^{\tilde{\mathfrak b}_+}=V_{\kappa_c}(\mathfrak g)^{\mathfrak g[\![t]\!]},\qquad
    (\operatorname{gr}V_{\kappa_c}(\mathfrak g))^{\tilde{\mathfrak b}_+}=(\operatorname{gr}V_{\kappa_c}(\mathfrak g))^{\mathfrak g[\![t]\!]}=\mathbb C[\overline P_{i,m}]_{i=1,\dots,\ell;\,m<0}.
\]

我们需要描述 $(\operatorname{gr}\mathbb M_{0,\kappa_c})^{\tilde{\mathfrak b}_+}$。首先注意
\[
    \operatorname{gr}\mathbb M_{0,\kappa_c}=\operatorname{Sym}\mathfrak g((t))/\mathfrak b_+\simeq\operatorname{Fun}\mathfrak g^*[\![t]\!]_{(-1)},
\]
其中
\[
    \mathfrak g^*[\![t]\!]_{(-1)}=((\mathfrak n_-)^*\otimes t^{-1})\oplus\mathfrak g^*[\![t]\!]\simeq(\mathfrak g((t))/\tilde{\mathfrak b}_+)^*,
\]
并且我们用 $\mathfrak g((t))$ 与 $\mathfrak g^*((t))$ 之间的配对 (8.1-1)。按此等同，映射 $\phi_{\mathrm{cl}}$ 成为由自然嵌入 $\mathfrak g^*[\![t]\!]\hookrightarrow\mathfrak g^*[\![t]\!]_{(-1)}$ 诱导的环同态
\[
    \operatorname{Fun}\mathfrak g^*[\![t]\!]_{(-1)}\longrightarrow\operatorname{Fun}\mathfrak g^*[\![t]\!].
\]

% [Tutor Note: 原书 p.223 页图在此式印 $\operatorname{gr}\mathbb M_{0,\kappa_c}=\operatorname{Sym}\mathfrak g((t))/\mathfrak b_+$（无波浪号），下一行对偶空间却印 $(\mathfrak g((t))/\tilde{\mathfrak b}_+)^*$（有波浪号）；原书两处确实不一致。译文照录页图，不擅改。]

设 $\mathfrak g$ 的基 $\{J^a\}$ 取得使其为两个子集之并，而这两个子集分别是 $\mathfrak b_+$ 与 $\mathfrak n_-$ 的基。令 $\overline J^a_n$ 为 $\hat{\mathfrak n}^{(-1)}_+$ 上由公式
\begin{equation}
    \overline J^a_n(A(t))=\operatorname{Res}_{t=0}\langle A(t),J^a\rangle t^n\,dt
    \tag{8.1-4}
\end{equation}
定义的多项式函数。则作为代数，$\operatorname{Fun}\mathfrak g^*[\![t]\!]_{(-1)}$ 由这些线性泛函 $\overline J^a_n$ 生成，其中 $n<0$，或 $n=0$ 且 $J^a\in\mathfrak n_-$。

% [Tutor Note: p.223 末“$n<0$，或 $n=0$ 且 $J^a\in\mathfrak n_-$”这一半开条件，是本节全部指标平移的来源：$\mathfrak g^*[\![t]\!]_{(-1)}$ 里允许 $t^{-1}$ 方向的分量，故 $\overline P_{i,m}$ 的指标比 $\operatorname{Inv}\mathfrak g^*[\![t]\!]$ 的情形整体上移 $d_i+1$。读后面 (8.1-3) 与 (8.1-5) 的指标差异时，务必回到这条约定。]

其次，我们照 §3.3.4 的办法，把生成函数
\[
    \overline J^a(z)=\sum_n\overline J^a_n z^{-n-1}
\]
（求和范围依 $J^a\in\mathfrak b_+$ 还是 $J^a\in\mathfrak n_-$ 分别取 $n<0$ 或 $n\le0$）代入 $\mathfrak g$ 上的不变多项式 $\overline P_i$（$i=1,\dots,\ell$），从而在 $\mathfrak g^*[\![t]\!]_{(-1)}$ 上构造 $\tilde{\mathfrak b}_+$-不变函数。但此时若 $J^a\in\mathfrak n_-$，$\overline J^a(z)$ 有非零的 $z^{-1}$ 系数，所得的级数
\[
    \overline P_i(\overline J^a(z))=\sum_{m\in\mathbb Z}\overline P_{i,m}z^{-m-1}
\]
对一切 $m<d_i$ 都有非零系数 $\overline P_{i,m}$。于是得到自然同态
\[
    \mathbb C[\overline P_{i,m_i}]_{i=1,\dots,\ell;\,m_i<d_i}\longrightarrow(\operatorname{Fun}\mathfrak g^*[\![t]\!]_{(-1)})^{\tilde{\mathfrak b}_+}.
\]

\begin{lemma}
    该同态是同构。
\end{lemma}

\begin{proof}
令
\[
    \mathfrak g^*[\![t]\!]_{(0)}=((\mathfrak n_-)^*\otimes1)\oplus(\mathfrak g^*\otimes t\mathbb C[[t]])=t\,\mathfrak g^*[\![t]\!]_{(-1)}.
\]
显然，$\mathfrak g^*[\![t]\!]_{(-1)}$ 与 $\mathfrak g^*[\![t]\!]_{(0)}$ 上 $\tilde{\mathfrak b}_+$-不变函数空间同构（尽管两者的 $\mathbb Z$-分次不同），故我们考虑后者。记 $\mathfrak g^*[\![t]\!]^{reg}_{(0)}$ 为 $\mathfrak g^*[\![t]\!]_{(0)}$ 与
\[
    J\mathfrak g^*_{reg}=\mathfrak g^*_{reg}\times(\mathfrak g^*\otimes t\mathbb C[[t]])
\]
之交。于是 $\mathfrak g^*[\![t]\!]^{reg}_{(0)}=((\mathfrak n_-)^{*,reg}\otimes1)\oplus(\mathfrak g^*\otimes t\mathbb C[[t]])$，其中 $(\mathfrak n_-)^{*,reg}=(\mathfrak n_-)^*\cap\mathfrak g^*_{reg}$ 是 $(\mathfrak n_-)^*$ 的稠密开子集，因而 $\hat{\mathfrak n}^{reg}_+$ 在 $\hat{\mathfrak n}_+$ 中稠密开。

回想定理 3.4.2 证明中引入的态射 $Jp:J\mathfrak g^*_{reg}\to J\mathcal P$。那里已证明，群 $JG=G[\![t]\!]$ 沿 $Jp$ 的纤维可迁地作用。设 $\widetilde B_+$ 是 $G[\![t]\!]$ 中对应李代数 $\tilde{\mathfrak b}_+\subset\mathfrak g[\![t]\!]$ 的子群。注意对任意 $x\in\mathfrak g^*[\![t]\!]^{reg}_{(0)}$，群 $\widetilde B_+$ 等于 $G[\![t]\!]$ 中所有满足 $g\cdot x\in\mathfrak g^*[\![t]\!]^{reg}_{(0)}$ 的元素 $g$ 构成的子群。因此 $\widetilde B_+$ 沿 $Jp$ 在 $\mathfrak g^*[\![t]\!]^{reg}_{(0)}$ 上的限制的纤维可迁地作用。

于是，与定理 3.4.2 的证明同样地，$\mathfrak g^*[\![t]\!]^{reg}_{(0)}$ 上 $\widetilde B_+$-不变（等价地，$\tilde{\mathfrak b}_+$-不变）多项式之环，就是 $Jp$ 下 $\mathfrak g^*[\![t]\!]^{reg}_{(0)}$ 在 $J\mathcal P$ 中的像上的函数之环。但由构造可知，$\mathfrak g^*[\![t]\!]^{reg}_{(0)}$ 在 $J\mathcal P$ 中的像，是由方程 $\overline P_{i,m}=0$（$i=1,\dots,\ell$；$m=-1$）确定的子空间。故 $\mathfrak g^*[\![t]\!]^{reg}_{(0)}$ 上 $\tilde{\mathfrak b}_+$-不变多项式之环等于 $\mathbb C[\overline P_{i,m_i}]_{i=1,\dots,\ell;\,m_i<-1}$。由于 $\mathfrak g^*[\![t]\!]^{reg}_{(0)}$ 在 $\mathfrak g^*[\![t]\!]_{(0)}$ 中稠密，这同样是 $\hat{\mathfrak n}_+$ 上的不变多项式之环。

当从 $\mathfrak g^*[\![t]\!]_{(0)}$ 过渡到 $\mathfrak g^*[\![t]\!]_{(-1)}$ 时，需计入不变多项式的指标平移 $\overline P_{i,m_i}\mapsto\overline P_{i,m_i+d_i+1}$，它对应平移 $\overline J^a_n\mapsto\overline J^a_{n+1}$。于是得到引理之断言。\transnote{「与定理 3.4.2 的证明同样地」省略了两步：其一，$\widetilde B_+$ 沿 $Jp$ 的纤维可迁地作用，其稳定化恰与 $\mathfrak g^*[\![t]\!]_{(0)}$ 中的负方向补空间一致；其二，像的切空间只由 $m=-1$ 的方程 $\overline P_{i,m}=0$ 切出，$m\ge0$ 的系数由等变性自动为零。缺这两步，由「纤维维数相等」推不出「像等于整条纤维」。落点：§8.1.1 引理 8.1.1 的证明，原书 p.224–225。[来源：本书 p.224–225；Frenkel–Ben-Zvi 2004, Ch. 18.]}
\end{proof}

% [Tutor Note: p.224–225 引理 8.1.1 的证明只写了“与定理 3.4.2 同样地”，省略的关键是两步：其一，$\widetilde B_+$ 的轨道恰是 $Jp$ 的纤维（因 $\widetilde B_+$ 在 $x$ 处的稳定化与 $\mathfrak g^*[\![t]\!]^{reg}_{(0)}$ 的“负方向补空间”一致）；其二，$\mathfrak g^*[\![t]\!]^{reg}_{(0)}$ 在 $J\mathcal P$ 中的像只由 $m=-1$ 那一族方程切出，因为 $m\ge0$ 的 $\overline P_{i,m}$ 自动为零（$\overline J^a(z)$ 只有 $z^{\ge0}$ 的幂）。细节见 [Frenkel and Ben-Zvi 2004] 第 18 章或 [Frenkel 2005b]。]

\begin{corollary}
    映射 $\phi_{\mathrm{cl}}$ 是满射。
\end{corollary}

\begin{proof}
映射 $\phi_{\mathrm{cl}}$ 相当于把 $\mathbb C[P_{i,m_i}]$（$i=1,\dots,\ell$；$m_i<d_i$）上的自由多项式代数，对由 $P_{i,m_i}$（$i=1,\dots,\ell$；$0\le m_i<d_i$）生成的理想取商。
\end{proof}

由构造，$\deg\overline P_{i,m}=m-d_i$。于是得到 $(\operatorname{gr}\mathbb M_{0,\kappa_c})^{\tilde{\mathfrak b}_+}=(\operatorname{Fun}\hat{\mathfrak n}^{(-1)}_+)^{\tilde{\mathfrak b}_+}$ 的如下特征标公式：
\[
    \operatorname{ch}(\operatorname{gr}\mathbb M_{0,\kappa_c})^{\tilde{\mathfrak b}_+}=\prod_{m>0}(1-q^m)^{-\ell}.
\]

回想由定理 6.3.3，Verma 模 $\mathbb M_{0,\kappa_c}$ 同构于 Wakimoto 模 $W^+_{0,\kappa_c}$。因此 $(\mathbb M_{0,\kappa_c})^{\tilde{\mathfrak b}_+}=(W^+_{0,\kappa_c})^{\tilde{\mathfrak b}_+}$。此外按引理 6.3.4，$(W^+_{0,\kappa_c})^{\tilde{\mathfrak b}_+}=\pi_0$，故其特征标也等于 $\prod_{m>0}(1-q^m)^{-\ell}$。因此自然嵌入
\[
    \operatorname{gr}\bigl(\mathbb M^{\tilde{\mathfrak b}_+}_{0,\kappa_c}\bigr)\hookrightarrow(\operatorname{gr}\mathbb M_{0,\kappa_c})^{\tilde{\mathfrak b}_+}
\]
是同构。

考虑交换图
\[
    \begin{CD}
    \operatorname{gr}\bigl(\mathbb M^{\tilde{\mathfrak b}_+}_{0,\kappa_c}\bigr) @>>> \operatorname{gr}\bigl(V_{\kappa_c}(\mathfrak g)^{\mathfrak g[\![t]\!]}\bigr)\\
    @VVV @VVV\\
    (\operatorname{gr}\mathbb M_{0,\kappa_c})^{\tilde{\mathfrak b}_+} @>>> (\operatorname{gr}V_{\kappa_c}(\mathfrak g))^{\mathfrak g[\![t]\!]}.
    \end{CD}
\]
由上面讨论，左竖箭头是同构。又由推论 8.1.2，下水平箭头是满射，故右竖箭头是满射。但按引理 3.3.1 它又是单射。于是得到同构
\[
    \operatorname{gr}\bigl(V_{\kappa_c}(\mathfrak g)^{\mathfrak g[\![t]\!]}\bigr)\simeq(\operatorname{gr}V_{\kappa_c}(\mathfrak g))^{\mathfrak g[\![t]\!]}.
\]
特别地，这意味着 $\operatorname{gr}\mathfrak z(\hat{\mathfrak g})=\operatorname{gr}(V_{\kappa_c}(\mathfrak g))^{\mathfrak g[\![t]\!]}$ 的特征标，等于由公式 (4.3-3) 给出的 $(\operatorname{gr}V_{\kappa_c}(\mathfrak g))^{\mathfrak g[\![t]\!]}$ 的特征标。由于 $\operatorname{ch}\mathfrak z(\hat{\mathfrak g})=\operatorname{ch}(\operatorname{gr}\mathfrak z(\hat{\mathfrak g}))$，我们得到
\begin{equation}
    \operatorname{ch}\mathfrak z(\hat{\mathfrak g})=\prod_{i=1}^{\ell}\prod_{n_i\ge d_i+1}(1-q^{n_i})^{-1}.
    \tag{8.1-5}
\end{equation}

于是得到如下结果（见 §4.3.1）：

\begin{theorem}
    中心 $\mathfrak z(\hat{\mathfrak g})$ “尽可能大”，即
    \[
        \operatorname{gr}\mathfrak z(\hat{\mathfrak g})=\operatorname{Inv}\mathfrak g^*[\![t]\!].
    \]
    于是存在中心元 $S_i\in\mathfrak z(\hat{\mathfrak g})\subset V_{\kappa_c}(\mathfrak g)$，其符号等于 $\overline P_{i,-1}\in\operatorname{Inv}\mathfrak g^*[\![t]\!]$（$i=1,\dots,\ell$），并且
    \[
        \mathfrak z(\hat{\mathfrak g})=\mathbb C[S_{i,(n)}]_{i=1,\dots,\ell;\,n<0}|0\rangle,
    \]
    其中 $S_{i,(n)}$ 是顶点算子 $Y(S_i,z)$ 的 Fourier 系数。
\end{theorem}

这是一个非平凡的结果，它告诉我们中心的许多结构信息。但我们并不满足，因为如 §3.5.2 所解释，我们想理解中心的几何意义，尤其想知道群 $\operatorname{Aut}\mathcal O$ 如何作用在 $\mathfrak z(\hat{\mathfrak g})$ 上。这一点体现在定理 4.3.2 中：它把 $\mathfrak z(\hat{\mathfrak g})$ 与 $\operatorname{Op}_{{}^LG}(D)$ 上的函数代数等同（注意定理 4.3.2 蕴含定理 8.1.3，见公式 (4.3-2)）。为证明它，我们需要再多费些力气，完成 §7.1 所述计划的第 5、6 步。\transnote{「尽可能大」回答的是分次不变量是否被取满：它断言 (8.1-2) 的单射实为同构，$\mathfrak g^*[\![t]\!]$ 的每个 $\mathfrak g[\![t]\!]$-不变多项式都来自 $\mathfrak z(\hat{\mathfrak g})$，符号映射不丢高阶修正项。与 (8.1-5) 的特征标逐项比较即知，$\mathfrak z(\hat{\mathfrak g})$ 由 $\ell$ 个 $S_i$ 及其全部 Fourier 分量自由生成。落点：§8.1.1 定理 8.1.3，原书 p.226。[来源：本书 p.226 与 (4.3-3)；定理 3.4.2.]}

% [Tutor Note: p.226 定理 8.1.3 的“尽可能大”是对 §4.3.1 中“中心是否取到分次伴随的全部不变量”这一问题的回答：单射 (8.1-2) 之所以可能不是满射，是因为 PBW 符号映射可能丢掉高阶修正；本定理说明它恰是同构，故 $\mathfrak z(\hat{\mathfrak g})$ 由 $\ell$ 个生成元 $S_i$ 及其全部 Fourier 分量自由生成，与 (8.1-5) 的特征标完全匹配。]

\subsection{中心与经典 W-代数}

按命题 7.3.6，$\mathfrak z(\hat{\mathfrak g})$ 含于 $\pi_0$ 上算子 $\overline V_i[1]$（$i=1,\dots,\ell$）的核之交中。现在计算这个交的特征标，并与 $\mathfrak z(\hat{\mathfrak g})$ 的特征标公式 (8.1-5) 比较，以证明 $\mathfrak z(\hat{\mathfrak g})$ 等于这些核之交。

首先，把该交与一族单参数顶点代数的经典极限等同起来，这族顶点代数称为 W-代数。给定单李代数 $\mathfrak g$ 及其上的不变内积 $\upsilon$，相应的 W-代数 $\mathcal W_\upsilon(\mathfrak g)$ 由量子 Drinfeld–Sokolov 约化在 [Feigin and Frenkel 1992] 中定义（亦见 [Frenkel and Ben-Zvi 2004] 第 15 章）。对 $\upsilon$ 的一般取值，顶点代数 $\mathcal W_\upsilon(\mathfrak g)$ 等于某个 Heisenberg 顶点代数中某些屏蔽算子的核之交。下面回忆这些算子的定义。\footnote{事实上，我们可以用这条性质对一切 $\upsilon$ 定义 $\mathcal W_\upsilon(\mathfrak g)$，见 [Feigin and Frenkel 1996]。}

考虑 §6.2.1 引入的 Heisenberg 李代数的另一个副本 $\hat{\mathfrak h}_\upsilon$。为避免混淆，把这个 Heisenberg 李代数的生成元记作
\[
    \mathbf b_{i,n},\qquad i=1,\dots,\ell;\ n\in\mathbb Z.
\]

我们有顶点算子
\[
    V^\upsilon_{-\alpha_i}(z):\pi^\upsilon_0\longrightarrow\pi^\upsilon_{-\alpha_i},
\]
它由公式 (7.2-13) 定义；令 $V^\upsilon_{-\alpha_i}[1]=\int V^\upsilon_{-\alpha_i}(z)\,dz$ 为其留数。我们称它为 W-代数屏蔽算子（以区别于前面定义的 Kac–Moody 代数屏蔽算子）。由于
\[
    V^\upsilon_{-\alpha_i}(z)=Y_{\pi^\upsilon_0,\pi^\upsilon_{-\alpha_i}}(|-\alpha_i\rangle,z)
\]
（记号见 §7.2.1），可知 $\pi^\upsilon_0$ 中 $V^\upsilon_{-\alpha_i}[1]$（$i=1,\dots,\ell$）的核之交是 $\pi^\upsilon_0$ 的顶点子代数。按 [Frenkel and Ben-Zvi 2004] 定理 15.4.12，对 $\upsilon$ 的一般取值，W-代数 $\mathcal W_\upsilon(\mathfrak g)$ 等于 $\pi^\upsilon_0$ 中算子 $V^\upsilon_{-\alpha_i}[1]$（$i=1,\dots,\ell$）的核之交。

我们关心 $\upsilon\to\infty$ 时 $\mathcal W_\upsilon(\mathfrak g)$ 的极限。为定义该极限，固定 $\mathfrak g$ 上的不变内积 $\upsilon_0$，记 $\epsilon$ 为 $\upsilon_0$ 与 $\upsilon$ 之比。利用 $\upsilon=\upsilon_0/\epsilon$ 诱导的 $\mathfrak h^*$ 与 $\mathfrak h$ 之间的等同，第 $i$ 个单根 $\alpha_i\in\mathfrak h^*$ 作为 $\mathfrak h$ 的元素有如下公式：
\[
    \alpha_i=\epsilon\frac{2}{\upsilon_0(h_i,h_i)}h_i.
\]
令
\begin{equation}
    \mathbf b'_{i,n}=\epsilon\frac{2}{\upsilon_0(h_i,h_i)}\mathbf b_{i,n},
    \tag{8.1-6}
\end{equation}
其中 $\mathbf b_{i,n}$ 是 $\hat{\mathfrak h}_\upsilon$ 的生成元。考虑 $\pi^\upsilon_0\otimes_{\mathbb C}\mathbb C[\epsilon]$ 中由 $\mathbf b'_{i,n}$（$i=1,\dots,\ell$；$n<0$）的一切单项式张成的 $\mathbb C[\epsilon]$-格。记 $\pi^\vee_0$ 为该格在 $\epsilon=0$ 处的特殊化；它是交换顶点代数。\transnote{这里的 $\epsilon$ 是内积之比（$\upsilon=\upsilon_0/\epsilon$，故 $\epsilon\to0$ 即 $\upsilon\to\infty$）；§8.1.4 讨论顶点 Poisson 结构时的 $\epsilon=(\kappa-\kappa_c)/\kappa_0$ 则是水平对临界水平的偏离。同一字母承担两个形变参数，分别见 (8.1-6) 与 (8.1-9)，读下文时须按上下文分清。落点：§8.1.2 与 §8.1.4，原书 p.227–229。[来源：本书 p.227–228.]}

在 $\epsilon\to0$ 的极限下，得到算子 $V^\upsilon_{-\alpha_i}[1]$ 的如下展开：
\[
    V^\upsilon_{-\alpha_i}[1]=\epsilon\frac{2}{\upsilon_0(h_i,h_i)}\mathbf V_i[1]+\cdots,
\]
其中省略号表示 $\epsilon$ 的高阶项；作用在 $\pi^\vee_0$ 上的算子 $\mathbf V_i[1]$\transnote{收缩算子写作 $\mathbf D_{\mathbf b'_{i,m-1}}$ 而非 $\mathbf D_{\mathbf b'_{i,m}}$，是为与 (8.1-8) 的 $\delta_{n,m}$ 配合：$\mathbf V_i[1]$ 取 $z^{-1}$ 的留数系数，作用时把指标推后一格，故 $\mathbf D$ 的下标要减一。若写成 $m$，(8.1-7) 与 (8.1-8) 的配对便不再自洽，(8.2-10) 处的递推也整体错位。落点：§8.1.2，(8.1-7)(8.1-8)，原书 p.227。[来源：本书 p.227 与 (7.2-13)；§7.2.1.]} 由公式
\begin{equation}
    \mathbf V_i[1]=\sum_{m\le0}\mathbf V_i[m]\mathbf D_{\mathbf b'_{i,m-1}}
    \tag{8.1-7}
\end{equation}
给出，其中
\begin{equation}
    \mathbf D_{\mathbf b'_{i,m}}\cdot\mathbf b'_{j,n}=a_{ij}\delta_{n,m},
    \tag{8.1-8}
\end{equation}
$(a_{ij})$ 是 $\mathfrak g$ 的 Cartan 矩阵，并且
\[
    \sum_{n\le0}\mathbf V_i[n]z^{-n}=\exp\Bigl(-\sum_{m>0}\frac{\mathbf b'_{i,-m}}{m}z^m\Bigr).
\]

% [Tutor Note: p.227 式 (8.1-6) 里 $\epsilon$ 的定义是 $\upsilon=\upsilon_0/\epsilon$，故 $\epsilon\to0$ 对应 $\upsilon\to\infty$；而 §8.1.4 中另用 $\kappa-\kappa_c=\epsilon\kappa_0$。同一个字母 $\epsilon$ 在两处分别是“内积之比”和“水平偏离临界水平之比”，这是全书最易混淆之处，读 (8.1-9) 与 (8.2-2) 附近时请随时回到各自的定义。]

% [Tutor Note: p.227 式 (8.1-7) 中 $\mathbf D_{\mathbf b'_{i,m-1}}$ 是与 $\mathbf b'_{i,m-1}$ 对偶的收缩算子（作用在 Fock 空间上），指标写成 $m-1$ 而与 (8.1-8) 的 $\delta_{n,m}$ 配合，等价于 $\mathbf V_i[1]$ 是 $z^{-1}$ 的留数系数；若把 $m-1$ 误写成 $m$，屏蔽算子的幂零性会失效。]

\subsection{Langlands 对偶李代数的出现}

算子 $\mathbf V_i[1]$（$i=1,\dots,\ell$）的核之交是 $\pi^\vee_0$ 的交换顶点子代数，记作 $\mathbf W(\mathfrak g)$，称为 $\mathfrak g$ 对应的经典 W-代数。

注意 $\pi^\vee_0$ 上的交换顶点代数结构与 $\upsilon_0$ 的选取无关。改变 $\upsilon_0$ 只使算子 $\mathbf V_i[1]$ 重新标度，故其核与 $\upsilon_0$ 无关。因此 $\mathbf W(\mathfrak g)$ 是 $\pi^\vee_0$ 的不依赖 $\upsilon_0$ 的交换顶点子代数。\footnote{然而下文将看到，$\pi^\vee_0$ 与 $\mathbf W(\mathfrak g)$ 还带有顶点 Poisson 代数结构，而这些结构确实依赖 $\upsilon_0$。}

比较公式 (7.3-13) 与 (8.1-7)，作代换 $b_{i,n}\mapsto-\mathbf b'_{i,n}$ 之后，算子 $\overline V_i[1]$ 变成算子 $\mathbf V_i[1]$，只是公式 (7.3-14) 中的矩阵系数 $a_{ji}$ 换成了公式 (8.1-8) 中的 $a_{ij}$。这不是笔误！这里有深刻原因：算子 $\overline V_i[1]$ 是 $\mathfrak g$ 的第 $i$ 个\emph{余根}对应的顶点算子的极限，而算子 $\mathbf V_i[1]$ 是 $\mathfrak g$ 的第 $i$ 个\emph{根}取负后对应的顶点算子的极限。\transnote{「这不是笔误」针对的是 $a_{ji}$ 变成 $a_{ij}$：代换 $b_{i,n}\mapsto-\mathbf b'_{i,n}$ 只改符号，真正把 Cartan 矩阵转置的是「$\mathfrak g$ 的第 $i$ 个余根」与「${}^L\mathfrak g$ 的第 $i$ 个根」的互换。两者合起来，$\pi_0$ 侧的 $\overline V_i[1]$ 才与 $\pi^\vee_0$ 侧的 $\mathbf V_i[1]$ 对上；这也是 Langlands 对偶李代数在此自然出现的地方。落点：§8.1.3，原书 p.228。[来源：本书 p.228；Feigin–Frenkel, \emph{Infinite analysis}, Adv. Ser. Math. Phys. 16, World Sci. 1992.]}

根与余根互换时，Cartan 矩阵转置。单李代数 $\mathfrak g$ 的转置 Cartan 矩阵是另一个单李代数的 Cartan 矩阵，即 $\mathfrak g$ 的 \emph{Langlands 对偶李代数}，记作 ${}^L\mathfrak g$。由于这个转置，可以把 ${}^L\mathfrak g$ 的 Cartan 子代数 ${}^L\mathfrak h$ 典范地等同于 $\mathfrak g$ 的 Cartan 子代数 $\mathfrak h$ 的对偶空间 $\mathfrak h^*$，于是 $\mathfrak g$ 的单根（$\mathfrak h^*$ 中的向量）成为 ${}^L\mathfrak g$ 的单余根（${}^L\mathfrak h$ 中的向量）。

把 $\mathfrak g$ 对应的 $\pi_0$ 与 ${}^L\mathfrak g$ 对应的 $\pi^\vee_0$ 通过 $b_{i,n}\mapsto-\mathbf b'_{i,n}$ 等同。于是附在 $\mathfrak g$ 上的算子 $\overline V_i[1]$ 变成附在 ${}^L\mathfrak g$ 上的算子 $\mathbf V_i[1]$。因此，单李代数 $\mathfrak g$ 对应的 $\pi_0$ 上算子 $\overline V_i[1]$（$i=1,\dots,\ell$）的核之交，同构于 ${}^L\mathfrak g$ 对应的 $\pi^\vee_0$ 上算子 $\mathbf V_i[1]$ 的核之交。

利用命题 7.3.6，现在得知 $\mathfrak z(\hat{\mathfrak g})$ 嵌入 $\pi^\vee_0$ 上算子 $\mathbf V_i[1]$（$i=1,\dots,\ell$）的核之交，即嵌入经典 W-代数 $\mathbf W({}^L\mathfrak g)$。此外有如下结果，其证明推迟到 §8.2.4。

\begin{lemma}
    $\mathbf W({}^L\mathfrak g)$ 的特征标，等于公式 (8.1-5) 给出的 $\mathfrak z(\hat{\mathfrak g})$ 的特征标。
\end{lemma}

% [Tutor Note: p.228 这段是全章的枢纽：$b_{i,n}\mapsto-\mathbf b'_{i,n}$ 中的负号与矩阵转置 $a_{ji}\to a_{ij}$ 合起来，把“$\mathfrak g$ 的余根”换成了“${}^L\mathfrak g$ 的根”，于是中心的对称代数（${}^L\mathfrak g$ 一侧）自然出现。若漏掉负号，$\mathbf V_i[1]$ 与 $\overline V_i[1]$ 的核将不再同构，定理 8.1.5 即不成立。]

于是得到如下结果。

\begin{theorem}
    作为交换顶点代数，中心 $\mathfrak z(\hat{\mathfrak g})$ 同构于 $\pi_0$ 上算子 $\overline V_i[1]$（$i=1,\dots,\ell$）的核之交，从而同构于经典 W-代数 $\mathbf W({}^L\mathfrak g)$。
\end{theorem}

这就完成了 §7.1 所述计划的第 5 步。

\subsection{顶点 Poisson 代数结构}

除交换顶点代数结构之外，$\mathfrak z(\hat{\mathfrak g})$ 与 $\mathbf W({}^L\mathfrak g)$ 还都带有\emph{顶点 Poisson 代数}结构；我们想证明定理 8.1.5 的同构与这些结构相容。

这里不给出顶点 Poisson 代数的精确定义，请读者参阅 [Frenkel and Ben-Zvi 2004] §16.2。我们只回忆：顶点 Poisson 代数 $P$ 特别是一个顶点李代数，于是可给它配上通常的李代数
\[
    \operatorname{Lie}(P)=P\otimes\mathbb C((t))/\operatorname{Im}(T\otimes1+1\otimes\partial_t)
\]
（见 [Frenkel and Ben-Zvi 2004] §16.1.7）。

按 [Frenkel and Ben-Zvi 2004] 命题 16.2.7，若 $V_\epsilon$ 是单参数顶点代数族，则 $V_0$ 的中心 $\mathcal Z(V_0)$ 自然获得顶点 Poisson 代数结构。具体说，在 $\epsilon=0$ 处，运算 $Y$ 限制在 $\mathcal Z(V_0)$ 上的极点部分消失，于是把 $Y$ 限制在 $\mathcal Z(V_0)$ 上的极点部分的 $\epsilon$-线性项定义为运算 $Y_-$。

固定 $\mathfrak g$ 上非零内积 $\kappa_0$，令 $\epsilon$ 为内积 $\kappa-\kappa_c$ 与 $\kappa_0$ 之比。把顶点代数 $V_\kappa(\mathfrak g)$ 视为以 $\kappa$ 为参数的单参数族。于是 $V_{\kappa_c}(\mathfrak g)$ 的中心 $\mathfrak z(\hat{\mathfrak g})$（对应 $\epsilon=0$）获得顶点 Poisson 结构。把这个带顶点 Poisson 结构的 $\mathfrak z(\hat{\mathfrak g})$ 记作 $\mathfrak z(\hat{\mathfrak g})_{\kappa_0}$。

其次，考虑 §6.2.1 引入的 Heisenberg 顶点代数 $\pi_0^{\kappa-\kappa_c}$。此后为避免混淆，我们写 $\pi_0^{\kappa-\kappa_c}(\mathfrak g)$ 以表明它附属于 $\mathfrak g$ 的 Cartan 子代数 $\mathfrak h$。仍令 $\epsilon$ 由公式 $\kappa-\kappa_c=\epsilon\kappa_0$ 定义。我们按如下方式把交换顶点代数 $\pi_0(\mathfrak g)$ 定义为 $\pi_0^{\kappa-\kappa_c}(\mathfrak g)$ 的 $\epsilon\to0$ 极限。回想 $\pi_0^{\kappa-\kappa_c}(\mathfrak g)$ 是以 $b_{i,n}$（$i=1,\dots,\ell$；$n\in\mathbb Z$）为生成元的 Heisenberg 李代数的 Fock 表示。其生成元的对易关系如下：
\begin{equation}
    [b_{i,n},b_{j,m}]=\epsilon n\kappa_0(h_i,h_j)\delta_{n,-m}.
    \tag{8.1-9}
\end{equation}
考虑 $\pi_0^{\kappa-\kappa_c}\otimes_{\mathbb C}\mathbb C[\epsilon]$ 中由 $b_{i,n}$（$i=1,\dots,\ell$；$n<0$）的一切单项式张成的 $\mathbb C[\epsilon]$-格。$\pi_0(\mathfrak g)$ 按定义就是该格在 $\epsilon=0$ 处的特殊化。这是交换顶点代数；但由于它是一族单参数顶点代数的极限，它获得顶点 Poisson 结构。该顶点 Poisson 结构由 Poisson 括号
\[
    \{b_{i,n},b_{j,m}\}=n\kappa_0(h_i,h_j)\delta_{n,-m}
\]
在 $\operatorname{Lie}(\pi_0(\mathfrak g))$ 中唯一确定，而这些括号可由公式 (8.1-9) 立即推出。特别地，该顶点 Poisson 结构依赖 $\kappa_0$，故我们写 $\pi_0(\mathfrak g)_{\kappa_0}$ 以表明这种依赖。\transnote{$\delta_{n,-m}$ 而非 $\delta_{n,m}$ 来自 Fock 空间的配对约定：$b_{i,n}$ 与 $b_{j,-m}$ 成对，故括号仅在 $n=-m$ 时非零；这与 (8.1-8) 的 $\mathbf D_{\mathbf b'_{i,m}}\cdot\mathbf b'_{j,n}=a_{ij}\delta_{n,m}$ 正好互相配合。取 $\epsilon\to0$ 后，$\frac1\epsilon[\, ,\,]$ 的极限即此括号，$\epsilon$ 充当普朗克常数。落点：§8.1.4，(8.1-9)，原书 p.229。[来源：本书 p.229；Frenkel–Ben-Zvi 2004, §16.2.]}

% [Tutor Note: p.229 式 (8.1-9) 中 $\delta_{n,-m}$（而不是 $\delta_{n,m}$）是 Fock 空间的标准约定：$b_{i,n}$ 与 $b_{j,-n}$ 配对，等价于 $\langle b_{i,n}b_{j,m}\rangle=n\kappa_0(h_i,h_j)\delta_{n,-m}$。取 $\epsilon\to0$ 后 $[\,,\,]$ 降为 Poisson 括号，$\epsilon$ 恰是“普朗克常数”，这正是“经典极限=顶点 Poisson 代数”的确切含义。]

回想同态 $w_{\kappa_c}:V_{\kappa_c}(\mathfrak g)\to W_{0,\kappa_c}=\mathbb M_{\mathfrak g}\otimes\pi_0(\mathfrak g)$ 可形变为同态 $w_\kappa:V_\kappa(\mathfrak g)\to W_{0,\kappa}=\mathbb M_{\mathfrak g}\otimes\pi_0^{\kappa-\kappa_c}(\mathfrak g)$。因此，$V_\kappa(\mathfrak g)$ 的运算 $Y$ 的极点部分限制在 $\mathfrak z(\hat{\mathfrak g})$ 上所得的 $\epsilon$-线性项（它用于定义 $\mathfrak z(\hat{\mathfrak g})$ 上的顶点 Poisson 结构），可通过把 $\pi_0^{\kappa-\kappa_c}(\mathfrak g)\subset W_{0,\kappa}$ 的运算 $Y$ 的极点部分的 $\epsilon$-线性项限制在 $\mathfrak z(\hat{\mathfrak g})\subset\pi_0(\mathfrak g)$ 上算出。而后者给 $\pi_0(\mathfrak g)$ 以顶点 Poisson 代数结构，记作 $\pi_0(\mathfrak g)_{\kappa_0}$。于是得到：

\begin{lemma}
    嵌入 $\mathfrak z(\hat{\mathfrak g})_{\kappa_0}\hookrightarrow\pi_0(\mathfrak g)_{\kappa_0}$ 是顶点 Poisson 代数的同态。相应局部李代数之间的映射 $\operatorname{Lie}(\mathfrak z(\hat{\mathfrak g})_{\kappa_0})\hookrightarrow\operatorname{Lie}(\pi_0(\mathfrak g)_{\kappa_0})$ 是李代数同态。
\end{lemma}

另一方面，考虑交换顶点代数 $\pi^\vee_0(\mathfrak g)$ 及其在 §8.1.2 定义的子代数 $\mathbf W(\mathfrak g)$。顶点代数 $\pi^\vee_0(\mathfrak g)$ 定义为一族顶点代数的极限，即 $\pi^\upsilon_0(\mathfrak g)$ 当 $\epsilon\to0$、$\upsilon=\upsilon_0/\epsilon$ 时的极限。因此 $\pi^\vee_0(\mathfrak g)$ 也带有顶点 Poisson 代数结构；把所得顶点 Poisson 代数记作 $\pi^\vee_0(\mathfrak g)_{\upsilon_0}$。其子代数 $\mathbf W(\mathfrak g)$ 是一族顶点子代数（即 $\mathcal W_\upsilon(\mathfrak g)$）的极限。因此 $\mathbf W(\mathfrak g)$ 是 $\pi^\vee_0(\mathfrak g)_{\upsilon_0}$ 的顶点 Poisson 子代数，记作 $\mathbf W(\mathfrak g)_{\upsilon_0}$。

于是我们看到，定理 8.1.5 的同构中出现的对象都带有自然的顶点 Poisson 代数结构。我们断言该同构实际上与这些结构相容。

下面更精确地说明。注意 $\mathfrak g$ 上非零不变内积 $\kappa_0$ 限制到 $\mathfrak h$ 给出 $\mathfrak h^*$ 上的非零内积，它是 ${}^L\mathfrak g$ 上不变内积 $\kappa^\vee_0$ 的限制。为避免混淆，把上述 Heisenberg 顶点代数的经典极限连同分别对应 $\kappa_0$ 与 $\kappa^\vee_0$ 的顶点 Poisson 结构，记作 $\pi_0(\mathfrak g)_{\kappa_0}$ 与 $\pi^\vee_0({}^L\mathfrak g)_{\kappa^\vee_0}$。有顶点 Poisson 代数同构
\[
    \iota:\pi_0(\mathfrak g)_{\kappa_0}\xrightarrow{\ \sim\ }\pi^\vee_0({}^L\mathfrak g)_{\kappa^\vee_0},
    \qquad b_{i,n}\mapsto-\mathbf b'_{i,n},
\]
其中 $\mathbf b'_{i,n}$ 由公式 (8.1-6) 给出。把同构 $\iota$ 限制到 $\pi_0(\mathfrak g)_{\kappa_0}$ 的子空间 $\bigcap_{1\le i\le\ell}\operatorname{Ker}\overline V_i[1]$ 上，得到顶点 Poisson 代数同构
\[
    \bigcap_{1\le i\le\ell}\operatorname{Ker}\overline V_i[1]\simeq\bigcap_{1\le i\le\ell}\operatorname{Ker}\mathbf V^\vee_i[1],
\]
其中 $\mathbf V^\vee_i[1]$ 表示附在 ${}^L\mathfrak g$ 上的算子 (8.1-7)。回想有如下顶点 Poisson 代数同构：
\[
    \mathfrak z(\hat{\mathfrak g})_{\kappa_0}\simeq\bigcap_{1\le i\le\ell}\operatorname{Ker}\overline V_i[1],
    \qquad
    \mathbf W({}^L\mathfrak g)_{\kappa^\vee_0}\simeq\bigcap_{1\le i\le\ell}\operatorname{Ker}\mathbf V^\vee_i[1].
\]
于是得到定理 8.1.5 的如下加强形式：

\begin{theorem}
    存在顶点 Poisson 代数的交换图
    \begin{equation}
    \begin{CD}
    \pi_0(\mathfrak g)_{\kappa_0} @>{\sim}>> \pi_0^\vee({}^L\mathfrak g)_{\kappa^\vee_0}\\
    @AAA @AAA\\
    \mathfrak z(\hat{\mathfrak g})_{\kappa_0} @>{\sim}>> \mathbf W({}^L\mathfrak g)_{\kappa^\vee_0}
    \end{CD}
    \tag{8.1-10}
    \end{equation}
\end{theorem}

\subsection{$\operatorname{Aut}\mathcal O$-模结构}

$\mathfrak z(\hat{\mathfrak g})_{\kappa_0}$ 与 $\mathbf W({}^L\mathfrak g)_{\kappa^\vee_0}$ 都带有群 $\operatorname{Aut}\mathcal O$ 的作用，我们断言定理 8.1.7 的同构与这两个作用相容。为看清这一点，我们把这两个作用都描述为来自顶点 Poisson 代数结构。

两种情形下，群 $\operatorname{Aut}\mathcal O$ 的作用都由李代数 $\operatorname{Der}_0\mathcal O\subset\operatorname{Der}\mathcal O$ 的作用取指数得到。就中心 $\mathfrak z(\hat{\mathfrak g})_{\kappa_0}$ 而言，$\operatorname{Der}\mathcal O$ 的作用是 $V_{\kappa_c}(\mathfrak g)$ 上自然作用的限制，而该自然作用来自 $\operatorname{Der}\mathcal O$ 通过无穷小变元替换在 $\hat{\mathfrak g}_{\kappa_c}$（保持其李子代数 $\mathfrak g[\![t]\!]$）上的作用。但在临界水平之外，即 $\kappa\ne\kappa_c$ 时，$\operatorname{Der}\mathcal O$ 的作用可通过 Virasoro 代数的作用得到，而后者来自公式 (6.2-5) 给出的共形向量 $\mathbf S_\kappa$。把该公式改写为
\[
    \mathbf S_\kappa=\frac{\kappa_0}{\kappa-\kappa_c}S_1,
\]
其中 $S_1$ 由公式 (3.1-1) 给出，$\kappa_0$ 是该公式中使用的内积。于是，当 $\kappa\ne\kappa_c$ 时，顶点算子
\[
    Y(\mathbf S_\kappa,z)=\sum_{n\in\mathbb Z}L_nz^{-n-2}
\]
的 Fourier 系数 $L_n$（$n\ge-1$）生成 $V_\kappa(\mathfrak g)$ 上的 $\operatorname{Der}\mathcal O$-作用。

在 $\kappa\to\kappa_c$ 的极限下，$\mathbf S_\kappa=\epsilon^{-1}S_1$，其中 $\epsilon=\frac{\kappa-\kappa_c}{\kappa_0}$ 如前。因此 $\operatorname{Der}\mathcal O$ 的作用通过 $\mathfrak z(\hat{\mathfrak g})_{\kappa_0}$ 上的顶点 Poisson 运算 $Y_-$ 得到：后者定义为当 $\epsilon\to0$ 时 $\epsilon^{-1}$ 乘以 $Y$ 的极点部分的极限（见 [Frenkel and Ben-Zvi 2004] §16.2），并作用在 $S_1\in\mathfrak z(\hat{\mathfrak g})_{\kappa_0}$ 上。

换言之，$\operatorname{Der}\mathcal O$-作用由级数
\[
    Y_-(S_1,z)=\sum_{n\ge-1}L_nz^{-n-2}
\]
的 Fourier 系数 $L_n$（$n\ge-1$）生成。于是我们看到，$\mathfrak z(\hat{\mathfrak g})_{\kappa_0}$ 上的自然 $\operatorname{Der}\mathcal O$-作用（从而 $\operatorname{Aut}\mathcal O$-作用），通过 $\mathfrak z(\hat{\mathfrak g})_{\kappa_0}$ 上的顶点 Poisson 代数结构，编码在向量 $S_1\in\mathfrak z(\hat{\mathfrak g})_{\kappa_0}$ 之中。注意该作用给 $\mathfrak z(\hat{\mathfrak g})_{\kappa_0}$ 以拟共形结构（见 §6.2.4）。

同样，$\mathbf W({}^L\mathfrak g)_{\kappa^\vee_0}$ 上有来自其顶点 Poisson 代数结构的 $\operatorname{Der}\mathcal O$-作用。生成该作用的向量也等于 W-代数 $\mathcal W_\upsilon({}^L\mathfrak g)$（它是共形顶点代数）中共形向量当 $\upsilon\to\infty$ 时的极限。该共形向量唯一，因为由 (8.1-5) 右端给出的 $\mathcal W_\upsilon({}^L\mathfrak g)$ 的特征标公式可见，$\mathcal W_\upsilon({}^L\mathfrak g)$ 中次数为 2 的齐次分量（一切共形向量都在其中）是一维的。\footnote{这把该向量确定到一个标量倍数，而标量由对易关系固定。}该向量当 $\upsilon\to\infty$ 的极限给出 $\mathbf W({}^L\mathfrak g)_{\kappa^\vee_0}$ 中的向量，记作 $\mathbf t$，使得 $Y_-(\mathbf t,z)$ 的 Fourier 系数生成 $\mathbf W({}^L\mathfrak g)_{\kappa^\vee_0}$ 上的 $\operatorname{Der}\mathcal O$-作用（下一节将解释该作用的几何意义）。由于这样的向量唯一，它必是 $S_1\in\mathfrak z(\hat{\mathfrak g})_{\kappa_0}$ 在同构 $\mathfrak z(\hat{\mathfrak g})_{\kappa_0}\simeq\mathbf W({}^L\mathfrak g)_{\kappa^\vee_0}$ 下的像。

于是我们断定：$\mathfrak z(\hat{\mathfrak g})_{\kappa_0}$ 与 $\mathbf W({}^L\mathfrak g)_{\kappa^\vee_0}$ 上的 $\operatorname{Der}\mathcal O$-作用（从而相应的 $\operatorname{Aut}\mathcal O$-作用），经由各自的顶点 Poisson 结构，由某些向量编码，而这些向量实际上等于共形向量的经典极限。在定理 8.1.7 的同构下，这些向量互相映射。于是得到：

\begin{proposition}
    交换图 (8.1-10) 与 $\operatorname{Der}\mathcal O$ 及 $\operatorname{Aut}\mathcal O$ 的作用相容。
\end{proposition}

当 $\kappa\ne\kappa_c$ 时，我们已在公式 (6.2-5) 中算出了 $\mathbf S_\kappa$ 在 $W_{0,\kappa}$ 中的像。取 $\kappa\to\kappa_c$ 的极限，得到像 $S_1=\epsilon\mathbf S_\kappa$ 属于 $\pi_0(\mathfrak g)_{\kappa_0}\subset W_{0,\kappa}$，且等于
\[
    \frac12\sum_{i=1}^{\ell}b_{i,-1}b^i_{-1}-\rho_{-2},
\]
其中 $\{b_i\}$ 与 $\{b^i\}$ 是关于 $S_1$ 的定义中所用内积 $\kappa_0$ 限制在 $\mathfrak h$ 上的一对对偶基，而 $\rho$ 是 $\mathfrak h^*$ 中 $\rho$ 在 $\kappa_0$ 诱导的同构 $\mathfrak h^*\simeq\mathfrak h$ 下对应的 $\mathfrak h$ 中元素。在定理 8.1.7 的同构下，该向量成为向量 $S_1\in\mathfrak z(\hat{\mathfrak g})_{\kappa_0}\subset\pi_0(\mathfrak g)_{\kappa_0}$，它负责其上的 $\operatorname{Der}\mathcal O$ 作用。\transnote{常数项 $-\rho_{-2}$ 来自正规序修正：$L_0$ 作用在二次的 Segal–Sugawara 向量上时产生一个常数平移，其值由 $\rho$ 的 $t^{-2}$ 系数给出。漏掉它，$\operatorname{Der}\mathcal O$ 的作用会多出常数中心项，与 (8.1-11) 的第二式不符。落点：§8.1.5，原书 p.232。[来源：本书 p.232 与 (6.2-5)、(3.1-1).]}

% [Tutor Note: p.232 这个公式 $\frac12\sum b_{i,-1}b^i_{-1}-\rho_{-2}$ 是 $S_1$ 在 Fock 空间中的像，注意第二项是常数 $\rho_{-2}$（即 $\rho$ 的 $t^{-2}$ 系数），它来自 Segal–Sugawara 向量的正规序修正；漏掉该项会使 $\operatorname{Der}\mathcal O$-作用多出一个常数中心项，与 (8.1-11) 的第二个公式不符。]

相应算子 $L_n\in\operatorname{Der}\mathcal O$（$n\ge-1$）在 $\pi_0(\mathfrak g)_{\kappa_0}$ 上的作用，由唯一确定的代数结构导子给出（见公式 (6.2-13)）。因此 $L_n$ 在 $\pi^\vee_0({}^L\mathfrak g)_{\kappa^\vee_0}$ 上的作用如下（回想在我们的同构 $\pi_0(\mathfrak g)_{\kappa_0}\simeq\pi^\vee_0({}^L\mathfrak g)_{\kappa^\vee_0}$ 下 $b_{i,n}\mapsto-\mathbf b'_{i,n}$）：
\begin{equation}
\begin{aligned}
    L_n\cdot\mathbf b'_{i,m}&=-m\mathbf b'_{i,n+m},&\qquad -1&\le n<-m,\\
    L_n\cdot\mathbf b'_{i,-n}&=-n(n+1),&\qquad n&>0,\\
    L_n\cdot\mathbf b'_{i,m}&=0,&\qquad n&>-m.
\end{aligned}
\tag{8.1-11}
\end{equation}
这些公式确定 $\pi^\vee_0({}^L\mathfrak g)_{\kappa^\vee_0}$ 上的 $\operatorname{Der}\mathcal O$-作用。由构造，$\mathbf W({}^L\mathfrak g)_{\kappa^\vee_0}\subset\pi^\vee_0({}^L\mathfrak g)_{\kappa^\vee_0}$ 在该作用下不变。\transnote{(8.1-11) 的三行按 $n$ 与 $-m$ 的大小分段：$n<-m$ 时平移，$n=-m$（此时 $m=-n<0$）时给出常数 $-n(n+1)$，$n>-m$ 时为零。第二行是唯一出现常数的地方，来自 $L_n$ 中正规序收缩项的残留；第一行限定 $n\ge-1$，因为 $\operatorname{Der}\mathcal O$ 只含 $n\ge-1$ 的模，不能简化成 $n<m$。落点：§8.1.5，(8.1-11)，原书 p.233。[来源：本书 p.233 与 (6.2-13).]}

特别注意，上述 $\operatorname{Der}\mathcal O$ 在 $\pi_0(\mathfrak g)_{\kappa_0}$ 与 $\pi^\vee_0({}^L\mathfrak g)_{\kappa^\vee_0}$ 上（从而在 $\mathfrak z(\hat{\mathfrak g})_{\kappa_0}$ 与 $\mathbf W({}^L\mathfrak g)_{\kappa^\vee_0}$ 上）的作用与内积 $\kappa_0$ 无关。

% [Tutor Note: p.233 式 (8.1-11) 的第二个公式左端下标是 $\mathbf b'_{i,-n}$，右端是 $-n(n+1)$，即 $L_n$ 作用在“负指标”生成元上给出常数；这与 $\mathbf b'_{i,n}$ 作为 $b$ 的线性组合、而 $L_n$ 是二次算子的结构有关。三行公式的条件分别读作 $-1\le n<-m$、$n>0$、$n>-m$，其中第一条不能简化为 $n<m$，因为 $\operatorname{Der}\mathcal O$ 只含 $n\ge-1$ 的模。]

\section{与 opers 上函数代数的等同}

\begin{keypoints}
    \item \textbf{核心思想：}经典 W-代数就是 oper 空间上的函数代数；Miura oper 空间就是 $H$-丛上联络空间，两者之间的 Miura 变换即遗忘映射。
    \item \textbf{数学工具：}第二个被联络保持的 Borel 约化（Miura 数据）、$w_0$-对偶 $H$-丛（引理 8.2.1）、$\mathfrak n$-不变的 Hamilton 作用与屏蔽算子。
    \item \textbf{关键结论：}命题 8.2.2（$\operatorname{MOp}_G(D)_{\mathrm{gen}}\simeq\operatorname{Conn}(\Omega^{\check\rho})_D$）、命题 8.2.3（$\operatorname{Fun}\operatorname{Op}_G(D)$ 是屏蔽算子核之交）、定理 8.2.5（$\mathfrak z(\hat{\mathfrak g})\simeq\operatorname{Fun}\operatorname{Op}_{{}^LG}(D)$）。
    \item \textbf{记号提醒：}原书在 §8.2 前半用 $\Omega^{\check\rho}$、后半用 $\Omega^{\rho}$ 与 $\Omega^{-\rho}$，三种写法经不变内积互相等同，译文照录；$p_{-1}=\sum_i f_i$ 是 oper 联络的“$-1$ 阶”部分。
\end{keypoints}

现在我们已经把中心 $\mathfrak z(\hat{\mathfrak g})_{\kappa_0}$ 与经典 W-代数 $\mathbf W({}^L\mathfrak g)_{\kappa^\vee_0}$ 等同起来，且这种等同与顶点 Poisson 代数结构及 $(\operatorname{Der}\mathcal O,\operatorname{Aut}\mathcal O)$-作用相容。证明定理 4.3.2 剩下的最后一步（§7.1 计划的第 6 步），是把 $\mathbf W({}^L\mathfrak g)_{\kappa^\vee_0}$ 与圆盘上 ${}^LG$-opers 空间上的函数代数 $\operatorname{Op}_{{}^LG}(D)$ 等同。本节完成这一步。

我们先引入 Miura opers，以及从一般 Miura opers 到 opers 的自然映射，称为 Miura 变换。然后证明：圆盘 $D$ 上 $G$-opers 与一般 Miura $G$-opers 的函数代数，分别同构于 $\mathbf W(\mathfrak g)_{\upsilon_0}$ 与 $\pi^\vee_0(\mathfrak g)_{\upsilon_0}$。此外，对应 Miura 变换的同态 $\mathbf W(\mathfrak g)_{\upsilon_0}\to\pi^\vee_0(\mathfrak g)_{\upsilon_0}$ 正是 §8.1.2 构造的嵌入。这将使我们能把中心 $\mathfrak z(\hat{\mathfrak g})$ 与 $D$ 上 ${}^LG$-opers 的函数代数等同。

% [Tutor Note: p.233 这段把 §8.1 的纯表示论结论翻译成几何命题：$\mathbf W(\mathfrak g)$ 不是抽象代数，而是 $\operatorname{Op}_G(D)$ 上的函数；屏蔽算子于是获得几何意义——它们是 $\mathfrak n((t))$ 的无穷小规范作用。这一步是整本书“中心 = opers 上函数”的核心机制。]

\subsection{Miura opers}

设 $G$ 是伴随型的单李群。

$X$（光滑曲线，或 $D$，或 $D^\times$）上的 \emph{Miura $G$-oper} 按定义是一个四元组 $(\mathcal F,\nabla,\mathcal F_{\mathcal B},\mathcal F'_{\mathcal B})$，其中 $(\mathcal F,\nabla,\mathcal F_{\mathcal B})$ 是 $X$ 上的 $G$-oper，而 $\mathcal F'_{\mathcal B}$ 是 $\mathcal F$ 的另一个被 $\nabla$ 保持的 $B$-约化。

考虑圆盘 $D$ 上 Miura $G$-opers 的空间 $\operatorname{MOp}_G(D)$。$\mathcal F$ 的被联络 $\nabla$ 保持的 $B$-约化，由 $\mathcal F$ 在原点 $0\in D$ 处纤维 $\mathcal F_0$ 的一个 $B$-约化唯一确定（回想一切 $G$-opers 的底层 $G$-丛彼此同构）。这样的约化之集，是旗簇 $G/B$ 的 $\mathcal F_0$-扭 $(G/B)_{\mathcal F_0}$。设 $\mathcal F_{\mathrm{univ}}$ 是 $\operatorname{Op}_G(D)$ 上的泛 $G$-丛，其在 oper $(\mathcal F,\nabla,\mathcal F_{\mathcal B})$ 处的纤维为 $\mathcal F_0$。于是得到
\[
    \operatorname{MOp}_G(D)\simeq(G/B)_{\mathcal F_{\mathrm{univ}}}=\mathcal F_{\mathrm{univ}}\times_G G/B.
\]

若两个 $B$-约化 $\mathcal F_{\mathcal B}$ 与 $\mathcal F'_{\mathcal B}$ 处于一般相对位置，则称 Miura $G$-oper 是\emph{一般的}（generic）。把 $D$ 上一般 Miura opers 的空间记作 $\operatorname{MOp}_G(D)_{\mathrm{gen}}$。

有自然的遗忘态射 $\operatorname{MOp}_G(D)_{\mathrm{gen}}\to\operatorname{Op}_G(D)$。群
\[
    N_{\mathcal F_{\mathcal B,0}}=\mathcal F_{\mathcal B,0}\times_B N,
\]
其中 $\mathcal F_{\mathcal B,0}$ 是 $\mathcal F_{\mathcal B}$ 在 $0$ 处的纤维，作用在 $(G/B)_{\mathcal F_0}$ 上，而一般约化的子集就是 $(G/B)_{\mathcal F_0}$ 的开 $N_{\mathcal F_{\mathcal B,0}}$-轨道。该轨道实际上是 $N_{\mathcal F_{\mathcal B,0}}$-torsor。因此，固定底层 $G$-oper $(\mathcal F,\nabla,\mathcal F_{\mathcal B})$ 的一般 Miura opers 空间是主 $N_{\mathcal F_{\mathcal B,0}}$-丛。\transnote{「一般的」（generic）指两个 Borel 约化处于一般相对位置，即 $\mathfrak b_{\mathcal F_{\mathcal B}}\cap\mathfrak b_{\mathcal F'_{\mathcal B}}\simeq\mathfrak h$；几何上等价于 $(G/B)_{\mathcal F_0}$ 落在开胞腔 $\mathcal U$ 中。因此 $\operatorname{MOp}_G(D)_{\mathrm{gen}}$ 是 $\operatorname{Op}_G(D)$ 上的主 $N$-丛，而 Miura 变换只在 gen 上有定义。落点：§8.2.1，原书 p.234 与 §4.2.4。[来源：本书 p.234；(8.2-1) 与 §4.2.4.]}

这可以改述如下。设 $\mathcal F_{\mathcal B,\mathrm{univ}}$ 是泛 $G$-丛 $\mathcal F_{\mathrm{univ}}$ 的 $B$-约化，其在 oper $(\mathcal F,\nabla,\mathcal F_{\mathcal B})$ 处的纤维为 $\mathcal F_{\mathcal B,0}$。则
\begin{equation}
    \operatorname{MOp}_G(D)_{\mathrm{gen}}=\mathcal F_{\mathcal B,\mathrm{univ}}\times_B\mathcal U,
    \tag{8.2-1}
\end{equation}
其中 $\mathcal U\simeq N$ 是 $G/B$ 中的开 $B$-轨道。

现在把 $\operatorname{MOp}_G(D)_{\mathrm{gen}}$ 与 $H$-联络的空间等同。对 $X$ 上的一般 Miura oper $(\mathcal F,\nabla,\mathcal F_{\mathcal B},\mathcal F'_{\mathcal B})$，考虑对应的 $H$-丛 $\mathcal F_H=\mathcal F_{\mathcal B}/N$ 与 $\mathcal F'_H=\mathcal F'_{\mathcal B}/N$。若 $\mathcal P$ 是 $H$-丛，用 $H$ 的对应于 $G$ 的外尔群最长元的自同构 $w_0$ 作用在它上面，得到新的 $H$-丛，记作 $w_0^*(\mathcal P)$。

% [Tutor Note: p.234 “一般的”指的是两个 Borel 约化处于一般相对位置，即 $\mathfrak b_{\mathcal F_{\mathcal B}}$ 与 $\mathfrak b_{\mathcal F'_{\mathcal B}}$ 的交为 $\mathfrak h$。这个条件在几何上等价于 $(G/B)_{\mathcal F_0}$ 落在开胞腔 $\mathcal U$ 中，因此 $\operatorname{MOp}_G(D)_{\mathrm{gen}}$ 是 $\operatorname{Op}_G(D)$ 上的主 $N$-丛，而不是整个 $\operatorname{MOp}_G(D)$。——这也解释了为什么后文的 Miura 变换只定义在 gen 上。]

\begin{lemma}
    对一般的 Miura oper $(\mathcal F,\nabla,\mathcal F_{\mathcal B},\mathcal F'_{\mathcal B})$，$H$-丛 $\mathcal F'_H$ 同构于 $w_0^*(\mathcal F_H)$。
\end{lemma}

\begin{proof}
考虑向量丛 $\mathfrak g_{\mathcal F}=\mathcal F\times_G\mathfrak g$，$\mathfrak b_{\mathcal F_{\mathcal B}}=\mathcal F_{\mathcal B}\times_B\mathfrak b$，$\mathfrak b_{\mathcal F'_{\mathcal B}}=\mathcal F'_{\mathcal B}\times_B\mathfrak b$。有处于一般位置的包含关系 $\mathfrak b_{\mathcal F_{\mathcal B}},\mathfrak b_{\mathcal F'_{\mathcal B}}\subset\mathfrak g_{\mathcal F}$。因此交 $\mathfrak b_{\mathcal F_{\mathcal B}}\cap\mathfrak b_{\mathcal F'_{\mathcal B}}$ 同构于 $\mathfrak b_{\mathcal F_{\mathcal B}}/[\mathfrak b_{\mathcal F_{\mathcal B}},\mathfrak b_{\mathcal F_{\mathcal B}}]$，即以 $\mathfrak h$ 为纤维的平凡向量丛。它自然作用在 $\mathfrak g_{\mathcal F}$ 上，在此作用下 $\mathfrak g_{\mathcal F}$ 分解为 $\mathfrak h$ 与线子丛 $\mathfrak g_{\mathcal F,\alpha}$（$\alpha\in\Delta$）的直和。此外 $\mathfrak b_{\mathcal F_{\mathcal B}}=\bigoplus_{\alpha\in\Delta_+}\mathfrak g_{\mathcal F,\alpha}$，$\mathfrak b_{\mathcal F'_{\mathcal B}}=\bigoplus_{\alpha\in\Delta_+}\mathfrak g_{\mathcal F,w_0(\alpha)}$。由于 $B$ 在 $\mathfrak n/[\mathfrak n,\mathfrak n]$ 上的作用经 $H=B/N$ 分解，我们得到
\[
    \mathcal F_H\times_H\bigoplus_{i=1}^{\ell}\mathbb C_{\alpha_i}\simeq\bigoplus_{i=1}^{\ell}\mathfrak g_{\mathcal F,\alpha_i},
    \qquad
    \mathcal F'_H\times_H\bigoplus_{i=1}^{\ell}\mathbb C_{\alpha_i}\simeq\bigoplus_{i=1}^{\ell}\mathfrak g_{\mathcal F,w_0(\alpha_i)}.
\]
于是得到
\[
    \mathcal F_H\times_H\mathbb C_{\alpha_i}\simeq\mathcal F'_H\times_H\mathbb C_{w_0(\alpha_i)},
    \qquad i=1,\dots,\ell.
\]
由于按假定 $G$ 是伴随型，上述相配线丛完全决定 $\mathcal F_H$ 与 $\mathcal F'_H$，而上述同构表明 $\mathcal F'_H\simeq w_0^*(\mathcal F_H)$。
\end{proof}

由于 $B$-丛 $\mathcal F'_{\mathcal B}$ 被 oper 联络 $\nabla$ 保持，我们得到 $\mathcal F'_H$ 上从而 $\mathcal F_H$ 上的联络 $\overline\nabla$。但按引理 4.2.1，$\mathcal F_H\simeq\Omega^{\check\rho}$。于是得到从 $D$ 上一般 Miura opers 空间 $\operatorname{MOp}_G(D)_{\mathrm{gen}}$ 到 $D$ 上 $H$-丛 $\Omega^{\check\rho}$ 的联络空间 $\operatorname{Conn}(\Omega^{\check\rho})_D$ 的态射 $\beta$。

具体说，$H$-丛 $\Omega^{\check\rho}$ 上的联络可用算子
\[
    \overline\nabla=\partial_t+\mathbf u(t),\qquad \mathbf u(t)\in\mathfrak h[\![t]\!]
\]
描述，其中 $t$ 是圆盘 $D$ 上的坐标，我们用它把 $\Omega^{\check\rho}$ 平凡化。设 $s$ 是新坐标，$t=\varphi(s)$。则同一联络将表示为 $\partial_s+\tilde{\mathbf u}(s)$，其中
\begin{equation}
    \tilde{\mathbf u}(s)=\varphi'\cdot\mathbf u(\varphi(s))-\check\rho\,\frac{\varphi''}{\varphi'}.
    \tag{8.2-2}
\end{equation}
该公式描述群 $\operatorname{Aut}\mathcal O$（及李代数 $\operatorname{Der}\mathcal O$）在 $\operatorname{Conn}(\Omega^{\check\rho})_D$ 上的作用。\transnote{$\check\rho\,\varphi''/\varphi'$ 是坐标变换下 $\Omega^{\check\rho}$ 上联络的修正项（射影联络型的 Schwarz 导数），其系数是基本余权之和 $\check\rho$，而非 $\rho$ 或 $\frac12\check\rho$。写成别的系数会与 (8.2-6) 的 $\operatorname{Der}\mathcal O$ 作用不符，Miura 变换也不再与 $\operatorname{Aut}\mathcal O$ 的作用相容。落点：§8.2.1，(8.2-2)，原书 p.235。[来源：本书 p.235；Beilinson–Drinfeld, \emph{Quantization of Hitchin's integrable system and Hecke eigensheaves}, 1997, §3.]}

\begin{proposition}
    映射 $\beta:\operatorname{MOp}_G(D)_{\mathrm{gen}}\to\operatorname{Conn}(\Omega^{\check\rho})_D$ 是同构。
\end{proposition}

\begin{proof}
我们构造反方向的映射 $\tau$。设给定 $D$ 上 $H$-丛 $\Omega^{\check\rho}$ 上的联络 $\overline\nabla$。按如下方式给它配一个一般 Miura oper。选取同态 $B\to H$ 的一个分裂 $H\to B$，并令 $\mathcal F=\Omega^{\check\rho}\times_H G$，$\mathcal F_{\mathcal B}=\Omega^{\check\rho}\times_H B$，其中 $G$ 与 $B$ 上的伴随作用取经由上述分裂得到的作用。分裂的选取还给出对面的 Borel 子群 $B_-$，它是与 $B$ 处于一般位置且含 $H$ 的唯一 Borel 子群。然后令 $\mathcal F'_{\mathcal B}=\Omega^{\check\rho}\times_H B_-w_0B$。

注意 $\mathcal F$ 上联络的空间同构于直积
\[
    \operatorname{Conn}(\Omega^{\check\rho})_D\times\bigoplus_{\alpha\in\Delta}\omega^{\alpha(\check\rho)+1}.
\]
对应负单根的子空间同构于空间 $\bigl(\bigoplus_{i=1}^{\ell}\mathfrak g_{-\alpha_i}\bigr)\otimes\mathcal O$。对每个 $i=1,\dots,\ell$ 取 $\mathfrak g_{-\alpha_i}$ 的一个基元 $f_i$，现在构造该空间的元素 $p_{-1}=\sum_{i=1}^{\ell}f_i$。令 $\nabla=\overline\nabla+p_{-1}$。由构造，$\nabla$ 与 $B$-约化 $\mathcal F_{\mathcal B}$ 处于正确的相对位置，并保持 $B$-约化 $\mathcal F'_{\mathcal B}$。因此四元组 $(\mathcal F,\nabla,\mathcal F_{\mathcal B},\mathcal F'_{\mathcal B})$ 是 $D$ 上的一般 Miura oper。令 $\tau(\overline\nabla)=(\mathcal F,\nabla,\mathcal F_{\mathcal B},\mathcal F'_{\mathcal B})$。

该映射与分裂 $H\to B$ 的选取以及生成元 $f_i$（$i=1,\dots,\ell$）的选取无关。确实，改变分裂 $H\to B$ 相当于用 $N$ 中元素共轭旧的分裂，这等价于对 $\nabla$ 施加该元素的规范变换，故不改变底层 Miura oper 结构。同样，重标度生成元 $f_i$ 可由 $H$ 中常元素的规范变换实现，这仍不改变 Miura oper 结构。由构造显然 $\beta$ 与 $\tau$ 互为逆同构。
\end{proof}

在命题 8.2.2 的同构下，自然遗忘态射 $\operatorname{MOp}_G(D)_{\mathrm{gen}}\to\operatorname{Op}_G(D)$ 成为映射
\begin{equation}
    \mu:\operatorname{Conn}(\Omega^{\check\rho})_D\longrightarrow\operatorname{Op}_G(D).
    \tag{8.2-3}
\end{equation}
我们称此映射为 \emph{Miura 变换}。

Miura 变换 (8.2-3) 给出相应函数环之间的同态 $\tilde\mu:\operatorname{Fun}\operatorname{Op}_G(D)\to\operatorname{Fun}\operatorname{Conn}(\Omega^{\check\rho})_D$。每个空间都有 $\operatorname{Der}\mathcal O$ 的作用，而 $\tilde\mu$ 是这两个环之间 $\operatorname{Der}\mathcal O$-等变的同态。下面我们先把 $\operatorname{Fun}\operatorname{Conn}(\Omega^{\check\rho})_D$ 与 $\pi^\vee_0(\mathfrak g)$ 等同，并把 $\tilde\mu$ 的像与 W-代数屏蔽算子核之交等同。这将给出 $\operatorname{Fun}\operatorname{Op}_G(D)$ 与 $\mathbf W(\mathfrak g)$ 的等同。

% [Tutor Note: p.235 式 (8.2-2) 中的 $\check\rho\frac{\varphi''}{\varphi'}$ 项是坐标变换下 $H$-丛 $\Omega^{\check\rho}$ 上联络的修正项（射影联络式的 Schwarz 导数）；系数恰为 $\check\rho$（基本余权之和），不是 $\rho$。若把它写成 $\frac12$ 或 $\rho$，$\operatorname{Aut}\mathcal O$-作用的公式 (8.2-6) 与 (8.2-2) 便不再一致，这是本节最易出错处。]

\subsection{Miura 变换的具体实现}

如 §4.2.4 所述，若在圆盘上取坐标 $t$，可把每个 oper 联络写成标准形
\[
    \partial_t+p_{-1}+\sum_{i=1}^{\ell}v_i(t)\cdot\mathbf c_i,\qquad v_i(t)\in\mathbb C[[t]]
\]
（见 §4.2.4），其中
\[
    v_i(t)=\sum_{n<0}v_{i,n}t^{-n-1}.
\]
于是
\[
    \operatorname{Fun}\operatorname{Op}_G(D)=\mathbb C[v_{i,n_i}]_{i=1,\dots,\ell;\,n_i<0}.
\]
若在 $\mathfrak b$ 中取 Cartan 子代数 $\mathfrak h$，则由命题 8.2.2，可把每个一般 Miura oper 表示为如下类型的联络算子：
\begin{equation}
    \partial_t+p_{-1}+\mathbf u(t),\qquad \mathbf u(t)\in\mathfrak h[\![t]\!].
    \tag{8.2-4}
\end{equation}
令 $u_i(t)=\alpha_i(\mathbf u(t))$，$i=1,\dots,\ell$，且
\[
    u_i(t)=\sum_{n<0}u_{i,n}t^{-n-1}.
\]
则
\[
    \operatorname{Fun}\operatorname{MOp}_G(D)_{\mathrm{gen}}=\operatorname{Fun}\operatorname{Conn}(\Omega^{\check\rho})_D=\mathbb C[u_{i,n}]_{i=1,\dots,\ell;\,n<0}.
\]
于是 Miura 变换给出同态
\begin{equation}
    \tilde\mu:\mathbb C[v_{i,n_i}]_{i=1,\dots,\ell;\,n_i<0}\longrightarrow\mathbb C[u_{i,n}]_{i=1,\dots,\ell;\,n<0}.
    \tag{8.2-5}
\end{equation}

\textbf{例。}我们计算 $\mathfrak g=\mathfrak{sl}_2$ 时的 Miura 变换 $\mu$。此时 oper 形如
\[
    \partial_t+\begin{pmatrix}0&v(t)\\1&0\end{pmatrix},
\]
而一般 Miura oper 形如
\[
    \partial_t+\begin{pmatrix}\frac12u(t)&0\\1&-\frac12u(t)\end{pmatrix}.
\]
为计算 $\mu$，需找出 $N[\![t]\!]$ 中元素，使相应的规范变换把 Miura oper 变成 oper 形式。我们算得
\[
    \begin{pmatrix}1&-\frac12u(t)\\0&1\end{pmatrix}
    \left(\partial_t+\begin{pmatrix}\frac12u(t)&0\\1&-\frac12u(t)\end{pmatrix}\right)
    \begin{pmatrix}1&\frac12u(t)\\0&1\end{pmatrix}
    =\partial_t+\begin{pmatrix}0&\frac14u(t)^2+\frac12\partial_tu(t)\\1&0\end{pmatrix}.
\]
于是得到
\[
    \mu(u(t))=v(t)=\frac14u(t)^2+\frac12\partial_tu(t),
\]
它也可写成
\[
    \partial_t^2-v(t)=\Bigl(\partial_t+\frac12u(t)\Bigr)\Bigl(\partial_t-\frac12u(t)\Bigr).
\]
最初正是 R. Miura 引入这个变换，作为把 KdV 族与 mKdV 族的流交织起来的映射。

当 $\mathfrak g=\mathfrak{sl}_2$ 时，opers 是射影联络，Miura opers 是仿射联络（见例如 [Frenkel and Ben-Zvi 2004] 第 9 章）。Miura 变换不是别的，正是从仿射联络到射影联络的自然映射。

% [Tutor Note: p.237 的 $\mathfrak{sl}_2$ 例子请自己重算一遍：规范变换用的矩阵 $\begin{pmatrix}1&-\frac12u\\0&1\end{pmatrix}$ 属于 $N[\![t]\!]$，其逆为 $\begin{pmatrix}1&\frac12u\\0&1\end{pmatrix}$；直接乘开即得 $\frac14u^2+\frac12\partial_tu$，中间 $\partial_tu$ 项的系数 $\frac12$ 来自两处各 $\frac12$ 之和，若漏掉矩阵求导项（$\partial_t$ 作用在右侧矩阵上）会得到 $\frac14u^2$，那是常见错误。]

\subsection{屏蔽算子}

由构造，Miura 变换 (8.2-3) 从而同态 (8.2-5) 是 $\operatorname{Der}\mathcal O$-等变的。$\operatorname{Der}\mathcal O$ 在 $\operatorname{Fun}\operatorname{Op}_G(D)$ 上的作用由公式 (4.2-5) 得到，在代数 $\operatorname{Fun}\operatorname{Conn}(\Omega^{\check\rho})_D$ 上的作用由公式 (8.2-2) 得到。它翻译成生成元 $L_n=-t^{n+1}\partial_t$（$n\ge-1$）在 $u_{i,m}$ 上的如下显式公式：
\begin{equation}
\begin{aligned}
    L_n\cdot u_{i,m}&=-mu_{i,n+m},&\qquad -1&\le n<-m,\\
    L_n\cdot u_{i,-n}&=-n(n+1),&\qquad n&>0,\\
    L_n\cdot u_{i,m}&=0,&\qquad n&>-m.
\end{aligned}
\tag{8.2-6}
\end{equation}

回想 $\operatorname{MOp}_G(D)_{\mathrm{gen}}$ 作为 $B$-torsor $\mathcal U\subset G/B$ 的 $\mathcal F_{\mathcal B,\mathrm{univ}}$-扭的实现 (8.2-1)。如 §4.2.4 所述，对一切 opers $(\mathcal F,\mathcal F_{\mathcal B},\nabla)$，$B$-torsor $\mathcal F_{\mathcal B,0}$ 彼此可以典范等同。于是群 $N_{\mathcal F_{\mathcal B,0}}$ 在映射 $\operatorname{MOp}_G(D)_{\mathrm{gen}}\to\operatorname{Op}_G(D)$ 的纤维上可迁地作用。按命题 8.2.2，有同构 $\operatorname{MOp}_G(D)_{\mathrm{gen}}\simeq\operatorname{Conn}(\Omega^{\check\rho})_D$。因此 $N_{\mathcal F_{\mathcal B,0}}$ 沿 Miura 变换 $\mu$ 的纤维可迁地作用。于是同态 $\tilde\mu$ 的像等于 $\operatorname{Fun}\operatorname{Conn}(\Omega^{\check\rho})_D$ 的 $N_{\mathcal F_{\mathcal B,0}}$-不变子空间，从而等于 $\operatorname{Fun}\operatorname{Conn}(\Omega^{\check\rho})_D$ 的 $\mathfrak n_{\mathcal F_{\mathcal B,0}}$-不变子空间。

在 $\mathfrak b$ 中固定 Cartan 子代数 $\mathfrak h$ 与 $\mathcal F_{\mathcal B,0}$ 的一个平凡化。用此平凡化把扭 $\mathfrak n_{\mathcal F_{\mathcal B,0}}$ 与 $\mathfrak n$ 等同。现在按 $\mathfrak h$ 在 $\mathfrak n$ 上的作用选取 $\mathfrak n$ 的生成元 $e_i$（$i=1,\dots,\ell$），使它们与前面选取的 $f_i$ 一起满足 $\mathfrak g$ 的标准关系。$\operatorname{Conn}(\Omega^{\check\rho})_D$ 上的 $N_{\mathcal F_{\mathcal B,0}}$-作用于是给出 $e_i$ 在 $\operatorname{Conn}(\Omega^{\check\rho})_D$ 上的无穷小作用。下面计算 $\operatorname{Fun}\operatorname{Conn}(\Omega^{\check\rho})_D$ 上相应的导子。

$e_i$ 的作用由无穷小规范变换
\begin{equation}
    \delta\mathbf u(t)=[x_i(t)\cdot e_i,\ \partial_t+p_{-1}+\mathbf u(t)]
    \tag{8.2-7}
\end{equation}
给出，其中 $x_i(t)\in\mathbb C[\![t]\!]$ 满足 $x_i(0)=1$，且 (8.2-7) 的右端属于 $\mathfrak h[\![t]\!]$。可以验证，这些条件唯一确定 $x_i(t)$。确实，(8.2-7) 的右端为
\[
    x_i(t)\cdot\check\alpha_i-u_i(t)x_i(t)\cdot e_i-\partial_tx_i(t)\cdot e_i.
\]
故它属于 $\mathfrak h[\![t]\!]$ 当且仅当\transnote{(8.2-7) 右端第一项 $x_i(t)\check\alpha_i$ 来自 $[x_i(t)e_i,p_{-1}]$：因 $p_{-1}=\sum_j f_j$，而 $[e_i,f_i]=h_i$ 在 $\mathfrak h^*$ 中对应余根 $\check\alpha_i$。这里用的是余根而非根；若混用，(8.2-10) 的 Cartan 矩阵会整体转置，命题 8.2.3 便不再与 (8.1-7) 吻合。落点：§8.2.3，(8.2-7)，原书 p.238–239。[来源：本书 p.238–239；(8.1-7) 与 (8.2-10).]}
\begin{equation}
    \partial_tx_i(t)=-u_i(t)x_i(t).
    \tag{8.2-8}
\end{equation}

若写
\[
    x_i(t)=\sum_{n\le0}x_{i,n}t^{-n}
\]
并代入公式 (8.2-8)，得到系数 $x_{i,n}$ 满足如下递推关系：
\[
    nx_{i,n}=\sum_{k+m=n;\,k<0;\,m\le0}u_{i,k}x_{i,m},\qquad n<0.
\]
由此算得
\begin{equation}
    \sum_{n\le0}x_{i,n}t^{-n}=\exp\Bigl(-\sum_{m>0}\frac{u_{i,-m}}{m}t^m\Bigr).
    \tag{8.2-9}
\end{equation}
现在得到
\[
    \delta u_j(t)=\alpha_j(\delta\mathbf u(t))=a_{ij}x_i(t),
\]
其中 $(a_{ij})$ 是 $\mathfrak g$ 的 Cartan 矩阵。换言之，算子 $e_i$ 通过导子
\begin{equation}
    \sum_{j=1}^{\ell}a_{ij}\sum_{n\ge0}x_{i,n}\frac{\partial}{\partial u_{j,-n-1}}
    \tag{8.2-10}
\end{equation}
作用在代数 $\operatorname{Fun}\operatorname{Conn}(\Omega^{\check\rho})_D=\mathbb C[u_{i,n}]$ 上，其中 $x_{i,n}$ 由公式 (8.2-9) 给出。

现在，$\operatorname{Fun}\operatorname{Op}_G(D)$ 在 Miura 映射 $\tilde\mu$ 下的像，是 $\operatorname{Conn}(\Omega^{\check\rho})_D$ 上 $\mathfrak n$-不变函数之代数。这些正是被 $\mathfrak n$ 的生成元 $e_i$（$i=1,\dots,\ell$，由公式 (8.2-10) 给出）零化的函数。于是得到 $\operatorname{Fun}\operatorname{Op}_G(D)$ 作为 $\operatorname{Fun}\operatorname{Conn}(\Omega^{\check\rho})_D$ 子代数的如下刻画。

\begin{proposition}
    在 Miura 映射 $\tilde\mu$ 下，$\operatorname{Fun}\operatorname{Op}_G(D)$ 在 $\operatorname{Fun}\operatorname{Conn}(\Omega^{\check\rho})_D$ 中的像，等于公式 (8.2-10)（$i=1,\dots,\ell$）所给算子的核之交。
\end{proposition}

% [Tutor Note: p.238–239 式 (8.2-7) 右端的第一项是 $x_i(t)\check\alpha_i$，它来自 $[x_ie_i,p_{-1}]=x_i[e_i,f_i]=x_ih_i$，而 $h_i$ 在 $\mathfrak h^*\simeq\mathfrak h$ 下对应 $\check\alpha_i$；注意这里用的是余根 $\check\alpha_i$ 而不是 $\alpha_i$。若把 $u_i=\alpha_i(\mathbf u)$ 与 $\check\alpha_i$ 的指标混用，(8.2-10) 中的 Cartan 矩阵 $a_{ij}$ 会整体转置，命题 8.2.3 就不再与 (8.1-7) 吻合。]

\subsection{回到 W-代数}

比较公式 (8.2-10) 与 (8.1-7)，我们发现：若在 W-代数屏蔽算子 $\mathbf V_i[1]$ 的公式 (8.1-7) 中把 $\mathbf b'_{i,n}$ 换成 $u_{i,n}$，就得到公式 (8.2-10)。因此算子 (8.2-10) 的核之交，等于算子 $\mathbf V_i[1]$（$i=1,\dots,\ell$）的核之交。但后者正是经典 W-代数 $\mathbf W(\mathfrak g)_{\upsilon_0}$。

于是得到如下交换图：
\begin{equation}
\begin{CD}
\pi^\vee_0(\mathfrak g)_{\upsilon_0} @>{\sim}>> \operatorname{Fun}\operatorname{Conn}(\Omega^{\check\rho})_D\\
@AAA @AAA\\
\mathbf W(\mathfrak g)_{\upsilon_0} @>{\sim}>> \operatorname{Fun}\operatorname{Op}_G(D)
\end{CD}
\tag{8.2-11}
\end{equation}
其中上箭头是代数同构，在生成元上由 $\mathbf b'_{i,n}\mapsto u_{i,n}$ 给出。

我们也可以计算 $\mathbf W(\mathfrak g)\simeq\operatorname{Fun}\operatorname{Op}_G(D)$ 的特征标，从而证明引理 8.1.4。由命题 8.2.2，$N$ 在 $\operatorname{Conn}(\Omega^{\check\rho})_D$ 上的作用是自由的。因此作为向量空间
\[
    \operatorname{Fun}\operatorname{Conn}(\Omega^{\check\rho})_D\simeq\operatorname{Fun}\operatorname{Op}_G(D)\otimes\operatorname{Fun}N.
\]
李代数 $\mathfrak n$ 的生成元 $e_i$ 的作用由算子 $\mathbf V_i[1]$ 给出，每个算子关于上面引入的 $\mathbb Z_+$-分次次数为 $-1$。因此根生成元 $e_\alpha\in\mathfrak n$（$\alpha\in\Delta_+$）作为次数 $-\langle\alpha,\check\rho\rangle$ 的算子作用。这意味着，若给 $\mathfrak n$ 配主分次的负分次（使 $\deg e_\alpha=\langle\alpha,\check\rho\rangle$），则 $\mathfrak n$ 的作用保持 $\operatorname{Fun}\operatorname{Conn}(\Omega^{\check\rho})_D$ 上的 $\mathbb Z$-分次。关于该分次，$\operatorname{Fun}N$ 的特征标等于
\[
    \prod_{\alpha\in\Delta_+}(1-q^{\langle\alpha,\check\rho\rangle})^{-1}=\prod_{i=1}^{\ell}\prod_{n_i=1}^{d_i}(1-q^{n_i})^{-1}.
\]
另一方面，$\operatorname{Fun}\operatorname{Conn}(\Omega^{\check\rho})_D$ 的特征标等于
\[
    \prod_{n>0}(1-q^n)^{-\ell}.
\]
因此 $\operatorname{Fun}\operatorname{Op}_G(D)$ 的特征标等于
\[
    \prod_{n>0}(1-q^n)^{-\ell}\prod_{i=1}^{\ell}\prod_{n_i=1}^{d_i}(1-q^{n_i})^{-1}.
\]
于是 $\operatorname{Fun}\operatorname{Op}_G(D)\simeq\mathbf W(\mathfrak g)_{\upsilon_0}$ 的特征标由公式 (8.1-5) 的右端给出。这就证明了引理 8.1.4。\transnote{此处「先算特征标、再断言同构」，特征标一步用的是按高度分次 $\deg e_\alpha=\langle\alpha,\check\rho\rangle$ 的计数：正根的次数应凑成 $1,\dots,d_i-1$（每个 $i$），而 (8.2-11) 右端的上限印作 $d_i$。与 (8.1-3)、(8.1-5) 中 $n_i\ge d_i+1$ 的起算点合看，三处指标不能同时成立；究竟哪一处是笔误，译者无法仅凭本书判定，故译文与公式均照录原书，存疑待审。落点：§8.2.4，(8.2-11)，原书 p.240，对照 p.223 与 p.226。[来源：本书 p.240；Kostant, \emph{Amer. J. Math.} 85 (1963), 327–404.]}

由上述同构还得到 $\operatorname{Fun}\operatorname{Op}_G(D)$ 上的顶点 Poisson 代数结构（该结构也可用 Drinfeld–Sokolov 约化定义，见下文）。如前，这些结构依赖内积 $\upsilon_0$ 的选取，我们有时会把它用作下标，以标明所考虑的 Poisson 结构。

% [Tutor Note: p.240 这段特征标计算是引理 8.1.4 的证明，也是全章唯一“先算特征标、再断言同构”的地方。按正文的分次 $\deg e_\alpha=\langle\alpha,\check\rho\rangle$（即高度分次），标准恒等式给出 $\prod_{\alpha\in\Delta_+}(1-q^{\langle\alpha,\check\rho\rangle})^{-1}=\prod_{i=1}^{\ell}\prod_{n_i=1}^{d_i-1}(1-q^{n_i})^{-1}$，而原书右端上限印作 $d_i$。若改用 $d_i-1$，则 $\operatorname{ch}\operatorname{Fun}\operatorname{Op}_G(D)=\prod_{i=1}^{\ell}\prod_{n\ge d_i}(1-q^n)^{-1}$，与 (8.1-5) 的 $n_i\ge d_i+1$ 相差一次整体平移。三处指标不能同时成立，译文照录原书，请数学审读裁定 (8.1-3)、(8.1-5) 与 p.240 何者需改。]

上图给出 W-代数屏蔽算子 $\mathbf V_i[1]$ 的几何解释：它们对应 $\mathfrak n$ 的生成元 $e_i$ 在 $\operatorname{Fun}\operatorname{Conn}(\Omega^{\check\rho})_D$ 上的作用。

比较公式 (8.2-6) 与 (8.1-11)，得知 (8.2-11) 中的上同构是 $\operatorname{Der}\mathcal O$-等变的。由于子空间 $\mathbf W(\mathfrak g)_{\upsilon_0}$ 与 $\operatorname{Fun}\operatorname{Op}_G(D)$ 在 $\operatorname{Der}\mathcal O$-作用下稳定，下同构也是 $\operatorname{Der}\mathcal O$-等变的。\footnote{我们指出，屏蔽算子作为 $\pi^\vee_0(\mathfrak g)$ 或 $\operatorname{Fun}\operatorname{Conn}(\Omega^{\check\rho})_D$ 的导子，与 $\operatorname{Der}\mathcal O$ 的作用不可交换。但若把它们看作从 $\pi^\vee_0(\mathfrak g)$ 到另一个模的算子（该模作为向量空间同构于 $\pi^\vee_0(\mathfrak g)$，但 $\operatorname{Der}\mathcal O$ 的作用被修改），则可使其与 $\operatorname{Der}\mathcal O$ 交换。此处不用这一事实，好奇的读者可参阅 [Feigin and Frenkel 1996]。}

此外，$\operatorname{Der}\mathcal O$ 在以上一切代数上的作用都与内积 $\upsilon_0$ 无关。于是得到：

\begin{theorem}
    图 (8.2-11) 与 $\operatorname{Der}\mathcal O$ 的作用相容。
\end{theorem}

\subsection{证明的完成}

现在把 $\mathfrak g$ 换成其 Langlands 对偶李代数 ${}^L\mathfrak g$。于是得到如下顶点 Poisson 代数的交换图
\begin{equation}
\begin{CD}
\pi^\vee_0({}^L\mathfrak g)_{\upsilon_0} @>{\sim}>> \operatorname{Fun}\operatorname{Conn}(\Omega^{\rho})_{D,\upsilon_0}\\
@AAA @AAA\\
\mathbf W({}^L\mathfrak g)_{\upsilon_0} @>{\sim}>> \operatorname{Fun}\operatorname{Op}_{{}^LG}(D)_{\upsilon_0}
\end{CD}
\tag{8.2-12}
\end{equation}
它与 $\operatorname{Der}\mathcal O$ 及 $\operatorname{Aut}\mathcal O$ 的作用相容。这就完成了 §7.1 计划的第六步，也是最后一步。我们现在已汇集了描述 $V_{\kappa_c}(\mathfrak g)$ 的中心 $\mathfrak z(\hat{\mathfrak g})$ 所需的全部部件。

把交换图 (8.2-12) 与 (8.1-10) 合并，并计入定理 8.2.4、定理 8.1.7 与命题 8.1.8，得到如下结果（此处 ${}^LG$ 指 ${}^L\mathfrak g$ 的内自同构群）。

\begin{theorem}
    存在同构 $\mathfrak z(\hat{\mathfrak g})_{\kappa_0}\simeq\operatorname{Fun}\operatorname{Op}_{{}^LG}(D)_{\kappa^\vee_0}$，它保持两边的顶点 Poisson 结构与 $\operatorname{Der}\mathcal O$-模结构。此外，它嵌入带 $\operatorname{Der}\mathcal O$-作用的顶点 Poisson 代数的交换图
    \begin{equation}
    \begin{CD}
    \pi_0(\mathfrak g)_{\kappa_0} @>{\sim}>> \operatorname{Fun}\operatorname{Conn}(\Omega^{\rho})_{D,\kappa^\vee_0}\\
    @AAA @AAA\\
    \mathfrak z(\hat{\mathfrak g})_{\kappa_0} @>{\sim}>> \operatorname{Fun}\operatorname{Op}_{{}^LG}(D)_{\kappa^\vee_0}
    \end{CD}
    \tag{8.2-13}
    \end{equation}
    其中上箭头在生成元上由 $b_{i,n}\mapsto-u_{i,n}$ 给出。
\end{theorem}

% [Tutor Note: p.241 (8.2-12) 与 (8.2-13) 的右上角写法不同：前者用 $\operatorname{Fun}\operatorname{Conn}(\Omega^\rho)_{D,\upsilon_0}$，后者用 $\operatorname{Fun}\operatorname{Conn}(\Omega^\rho)_{D,\kappa^\vee_0}$。这反映了两条路线——一条经 ${}^L\mathfrak g$ 的 Heisenberg 顶点代数（内积 $\upsilon_0$），一条经 $\hat{\mathfrak g}$ 的水平形变（内积 $\kappa_0$）；两者由 $\mathfrak g\simeq{}^L\mathfrak g$ 的内积等同而一致。读图时不要把它们当作不同的对象。]

特别地，我们现已证明了定理 4.3.2。

注意定理 8.2.5 与前面 §6.2.5 中给出的 $\pi_0(\mathfrak g)$ 的描述一致。按该描述，$b_i(t)$ 在 $\operatorname{Der}\mathcal O$ 作用下按圆盘 $D=\operatorname{Spec}\mathbb C[\![t]\!]$ 上 ${}^LH$-丛 $\Omega^{-\rho}$ 上联络的分量变换。因此 $\pi_0(\mathfrak g)$ 以 $\operatorname{Der}\mathcal O$-等变的方式，等同于 $\Omega^{-\rho}$ 上联络空间上的函数代数 $\operatorname{Fun}\operatorname{Conn}(\Omega^{-\rho})_D$。于是图 (8.2-13) 中的上同构可写成同构
\begin{equation}
    \operatorname{Fun}\operatorname{Conn}(\Omega^{-\rho})_D\simeq\operatorname{Fun}\operatorname{Conn}(\Omega^{\rho})_D.
    \tag{8.2-14}
\end{equation}
在此同构下，$\Omega^{\rho}$ 上的联络 $\partial_t+\mathbf u(t)$ 映为对偶 ${}^LH$-丛 $\Omega^{-\rho}$ 上的对偶联络 $\partial_t-\mathbf u(t)$。这恰好对应定理 8.2.5 中两个代数生成元之间的映射 $b_{i,n}\mapsto-u_{i,n}$。

给定任意圆盘 $D_x$，可考虑我们的代数由 $\operatorname{Aut}\mathcal O$-torsor $\mathcal Aut_x$ 所作的扭（定义见 §6.2.6）。我们用下标 $x$ 标记这些扭。于是得到
\[
    \mathbf W(\mathfrak g)_{\upsilon_0,x}\simeq\operatorname{Fun}\operatorname{Op}(D_x)_{\upsilon_0},
    \qquad
    \mathfrak z(\hat{\mathfrak g})_{\kappa_0,x}\simeq\operatorname{Fun}\operatorname{Op}_{{}^LG}(D_x)_{\kappa^\vee_0}.
\]

\subsection{分次伴随代数}

本节在分次伴随空间的层面描述定理 8.2.5 的同构 $\mathfrak z(\hat{\mathfrak g})\simeq\operatorname{Fun}\operatorname{Op}_{{}^LG}(D)$。先描述 $\mathfrak z(\hat{\mathfrak g})$ 与 $\operatorname{Fun}\operatorname{Op}_{{}^LG}(D)$ 上的滤过。

$\mathfrak z(\hat{\mathfrak g})$ 上的滤过由泛包络代数 $U(\hat{\mathfrak g}_{\kappa_c})$ 上的 Poincaré–Birkhoff–Witt 滤过诱导，见 §3.3.3。由定理 8.1.5，$\operatorname{gr}\mathfrak z(\hat{\mathfrak g})=(\operatorname{gr}V_{\kappa_c}(\mathfrak g))^{\mathfrak g[\![t]\!]}$。而
\[
    \operatorname{gr}V_{\kappa_c}(\mathfrak g)=\operatorname{Sym}\mathfrak g((t))/\mathfrak g[\![t]\!]\simeq\operatorname{Fun}\mathfrak g^*[\![t]\!]dt,
\]
与坐标 $t$ 及 $\mathfrak g$ 上内积的选取无关。命题 3.4.2 给出了 $(\operatorname{gr}V_{\kappa_c}(\mathfrak g))^{\mathfrak g[\![t]\!]}$ 的描述。该描述的不依赖坐标的版本如下。令 $C_{\mathfrak g}=\operatorname{Spec}(\operatorname{Fun}\mathfrak g^*)^G$ 及
\[
    C_{\mathfrak g,\Omega}=\Omega\times_{\mathbb C^\times}C_{\mathfrak g},
\]
其中 $\Omega=\mathbb C[\![t]\!]dt$ 是 $D=\operatorname{Spec}\mathbb C[\![t]\!]$ 上的拓扑微分模。由命题 3.4.2 有典范且不依赖坐标的同构
\begin{equation}
    \operatorname{gr}\mathfrak z(\hat{\mathfrak g})\simeq\operatorname{Fun}C_{\mathfrak g,\Omega}.
    \tag{8.2-15}
\end{equation}

注意，选取 $(\operatorname{Fun}\mathfrak g^*)^G$ 的齐次生成元 $\overline P_i$（$i=1,\dots,\ell$）给出等同
\[
    C_{\mathfrak g,\Omega}\simeq\bigoplus_{i=1}^{\ell}\Omega^{\otimes(d_i+1)},
\]
但我们宁愿不用它，因为这样的生成元没有自然选取。

其次考虑映射
\[
    a:\mathfrak z(\hat{\mathfrak g})\longrightarrow\pi_0(\mathfrak g),
\]
它等于嵌入 $V_{\kappa_c}(\mathfrak g)\hookrightarrow W_{0,\kappa_c}$ 限制在 $\mathfrak z(\hat{\mathfrak g})$ 上。在命题 7.1.1 的证明中，我们描述了 $W_{0,\kappa_c}$ 上与 $V_{\kappa_c}(\mathfrak g)$ 的 PBW 滤过相容的滤过。这意味着映射 $\mathfrak z(\hat{\mathfrak g})\to\pi_0(\mathfrak g)$ 也与滤过相容。按 §6.2.5 的结果，有典范等同
\[
    \pi_0(\mathfrak g)=\operatorname{Fun}\operatorname{Conn}(\Omega^{-\rho})_D.
\]
${}^LH$-丛 $\Omega^{-\rho}$ 上联络的空间 $\operatorname{Conn}(\Omega^{-\rho})_D$ 是向量空间 ${}^L\mathfrak h\otimes\Omega=\mathfrak h^*\otimes\Omega$ 上的仿射空间。因此
\begin{equation}
    \operatorname{gr}\pi_0(\mathfrak g)=\operatorname{Fun}\mathfrak h^*\otimes\Omega.
    \tag{8.2-16}
\end{equation}
有 Harish--Chandra 同构 $(\operatorname{Fun}\mathfrak g^*)^G\simeq(\operatorname{Fun}\mathfrak h^*)^W$，其中 $W$ 是 $\mathfrak g$ 的外尔群，于是有嵌入 $(\operatorname{Fun}\mathfrak g^*)^G\hookrightarrow\operatorname{Fun}\mathfrak h^*$。该嵌入给出嵌入
\[
    \operatorname{Fun}C_{\mathfrak g,\Omega}\longrightarrow\operatorname{Fun}\mathfrak h^*\otimes\Omega.
\]
由命题 7.1.1 的证明可知，在等同 (8.2-15) 与 (8.2-16) 下，这正是映射
\[
    \operatorname{gr}a:\operatorname{gr}\mathfrak z(\hat{\mathfrak g})\longrightarrow\operatorname{gr}\pi_0(\mathfrak g).
\]
这样我们就描述了交换图 (8.2-13) 中左竖映射的分次伴随。

下面描述交换图 (8.2-13) 中右竖映射的分次伴随。令
\[
    C^\vee_{\mathfrak g}=\mathfrak g/G=\operatorname{Spec}(\operatorname{Fun}\mathfrak g^*)^G
\]
及
\[
    C^\vee_{\mathfrak g,\Omega}=\Omega\times_{\mathbb C^\times}C^\vee_{\mathfrak g}.
\]
遵循 [Beilinson and Drinfeld 1997] §3.1.14，我们把 $\operatorname{Op}_G(D)$ 与以 $C^\vee_{\mathfrak g,\Omega}$ 为模型的仿射空间等同。换言之，$\operatorname{Fun}\operatorname{Op}_G(D)$ 上有自然滤过，使得
\begin{equation}
    \operatorname{gr}\operatorname{Op}_G(D)=\operatorname{Fun}C^\vee_{\mathfrak g,\Omega}.
    \tag{8.2-17}
\end{equation}
构造这样的滤过，等同于构造一个平坦的 $\mathbb C[h]$-代数，使其在 $h=1$ 处的特殊化为 $\operatorname{Fun}\operatorname{Op}_G(D)$，在 $h=0$ 处的特殊化为 $\operatorname{Fun}C^\vee_{\mathfrak g,\Omega}$。考虑 $D$ 上 $h$-opers 空间 $\operatorname{Op}_{G,h}(X)$ 上的函数代数。$h$-oper 的定义与 oper 相同，只是把形如 (4.2-1) 的联络换成 $h$-联络
\[
    h\partial_t+\sum_{i=1}^{\ell}\psi_i(t)f_i+\mathbf v(t).
\]
可以证明该代数在 $\mathbb C[h]$ 上平坦，办法与引理 4.2.2 的证明完全一样：证明每个 $h$-oper 都有标准形
\[
    h\partial_t+p_{-1}+\mathbf v(t),\qquad \mathbf v(t)\in V_{\mathrm{can}}((t)).
\]
显然 $\operatorname{Op}_{G,1}(D)=\operatorname{Op}_G(D)$。为看出 $\operatorname{Op}_{G,0}(D)=C^\vee_{\mathfrak g,\Omega}$，注意由定义 $\operatorname{Op}_{G,0}(D)=\Gamma(D,\Omega\times_{\mathbb C^\times}V_{\mathrm{can}})$，其中 $\mathbb C^\times$ 在 $V_{\mathrm{can}}$ 上的作用为 $a\mapsto a\operatorname{Ad}_{\rho(a)}$。但由 [Kostant 1963] 有 $\mathbb C^\times$-空间的典范同构 $V_{\mathrm{can}}\simeq\mathfrak g/G$。这意味着 $\operatorname{Op}_{G,0}(D)=C^\vee_{\mathfrak g,\Omega}$，从而 (8.2-17)。

% [Tutor Note: p.243–244 的 $h$-oper 是“把 oper 联络整体乘 $h$”所得的形变；$h=1$ 回到 oper，$h=0$ 退化为 $V_{\mathrm{can}}$ 的截面。平坦性（即 $\operatorname{Fun}\operatorname{Op}_{G,h}(D)$ 无挠）靠“每个 $h$-oper 有标准形”这一逐项计算，与引理 4.2.2 完全平行；这一步是把 $\operatorname{Op}_G(D)$ 的滤过与 $C^\vee_{\mathfrak g,\Omega}$ 的仿射结构对上的关键，不能跳过。]

现在考虑 ${}^LH$-丛 $\Omega^{\rho}$ 上的联络空间 $\operatorname{Conn}(\Omega^{\rho})_D$。它是向量空间 ${}^L\mathfrak h\otimes\Omega$ 上的仿射空间。因此
\begin{equation}
    \operatorname{gr}\operatorname{Fun}\operatorname{Conn}(\Omega^{\rho})_D=\operatorname{Fun}{}^L\mathfrak h\otimes\Omega.
    \tag{8.2-18}
\end{equation}
有 Harish--Chandra 同构 $(\operatorname{Fun}\mathfrak g)^G\simeq(\operatorname{Fun}\mathfrak h)^W$，其中 $W$ 是 $\mathfrak g$ 的外尔群，于是有嵌入 $(\operatorname{Fun}\mathfrak g)^G\hookrightarrow\operatorname{Fun}\mathfrak h$。该嵌入给出嵌入
\[
    \operatorname{Fun}C^\vee_{\mathfrak g,\Omega}\longrightarrow\operatorname{Fun}\mathfrak h\otimes\Omega.
\]
命题 8.2.2 证明中给出的 Miura 变换 $\operatorname{Conn}(\Omega^{\rho})_D\to\operatorname{Op}_{{}^LG}(D)$ 的显式构造表明，映射
\[
    \operatorname{Fun}\operatorname{Op}_{{}^LG}(D)\longrightarrow\operatorname{Fun}\operatorname{Conn}(\Omega^{\rho})_D
\]
保持滤过。此外，关于等同 (8.2-17) 与 (8.2-18)，其分次伴随映射不是别的，正是同态
\[
    \operatorname{Fun}C^\vee_{{}^L\mathfrak g,\Omega}\longrightarrow\operatorname{Fun}{}^L\mathfrak h\otimes\Omega.
\]
这描述了交换图 (8.2-13) 中右竖映射的分次伴随。

其次，上水平同构的分次伴随由自然复合
\[
    \operatorname{gr}\pi_0(\mathfrak g)\simeq\operatorname{Fun}\mathfrak h^*\otimes\Omega=\operatorname{Fun}{}^L\mathfrak h\otimes\Omega\simeq\operatorname{gr}\operatorname{Fun}\operatorname{Conn}(\Omega^{\rho})_D
\]
给出，再乘以算子 $(-1)^{\deg}$；该算子在次数 $n$ 的元素上取值 $(-1)^n$（这源于定理 8.2.5 中的映射 $b_{i,n}\mapsto-u_{i,n}$）。

注意有典范同构
\[
    (\operatorname{Fun}\mathfrak g^*)^G=(\operatorname{Fun}\mathfrak h^*)^W=(\operatorname{Fun}{}^L\mathfrak h)^W=(\operatorname{Fun}{}^L\mathfrak g)^{{}^LG},
\]
它们给出典范等同
\[
    \operatorname{Fun}C_{\mathfrak g,\Omega}=\operatorname{Fun}C^\vee_{{}^L\mathfrak g,\Omega}.
\]
现在，图 (8.2-13) 中一切映射都保持滤过。图的交换性意味着下水平同构
\begin{equation}
    \mathfrak z(\hat{\mathfrak g})\simeq\operatorname{Fun}\operatorname{Op}_{{}^LG}(D)
    \tag{8.2-19}
\end{equation}
的分次伴随，等于复合
\[
    \operatorname{gr}\mathfrak z(\hat{\mathfrak g})\simeq\operatorname{Fun}C_{\mathfrak g,\Omega}=\operatorname{Fun}C^\vee_{{}^L\mathfrak g,\Omega}\simeq\operatorname{gr}\operatorname{Fun}\operatorname{Op}_{{}^LG}(D)
\]
乘以算子 $(-1)^{\deg}$，它在次数 $n$ 的元素上取值 $(-1)^n$。于是得到：

\begin{theorem}
    同构 (8.2-19) 保持滤过。相应的分次伴随代数都同构于 $\operatorname{Fun}C_{\mathfrak g,\Omega}$。相应的分次伴随代数同构等于 $(-1)^{\deg}$。
\end{theorem}

\section{完备泛包络代数中心的描述}

\begin{keypoints}
    \item \textbf{核心思想：}把圆盘上的顶点 Poisson 结论“取包络代数”$\widetilde U$，即得完备泛包络代数中心 $Z(\hat{\mathfrak g})$ 与去心圆盘上 opers 函数代数之间的 Poisson 同构（定理 8.3.1，Drinfeld 猜想）。
    \item \textbf{数学工具：}$\widetilde U$ 函子把顶点 Poisson 代数变为拓扑 Poisson 代数；Drinfeld–Sokolov 约化（Hamilton 约化）给出 $\operatorname{Fun}\operatorname{Op}_G(D^\times)$ 上的同一 Poisson 结构。
    \item \textbf{关键结论：}定理 8.3.1（$Z(\hat{\mathfrak g})\simeq\operatorname{Fun}\operatorname{Op}_{{}^LG}(D^\times)_{\kappa^\vee_0}$）与定理 8.3.3（仿射 Harish--Chandra 同态 = Miura 变换）。
    \item \textbf{记号提醒：}$\Omega^{-\rho}$ 与 $\Omega^{\rho}$ 是互为对偶的 ${}^LH$-丛；Wakimoto 模 $W_{\overline\nabla}$ 对应点 $\overline\nabla\in\operatorname{Conn}(\Omega^{-\rho})_{D^\times}$。
\end{keypoints}

回想 §2.1.2 定义的 $\hat{\mathfrak g}_\kappa$ 的泛包络代数的完备化 $\widetilde U_\kappa(\hat{\mathfrak g})$。设 $Z(\hat{\mathfrak g})$ 为其中心。在定理 4.3.6 中，我们从定理 4.3.2（现已证明，见定理 8.2.5）导出了如下 $\operatorname{Aut}\mathcal O$-等变同构：
\begin{equation}
    Z(\hat{\mathfrak g})\simeq\operatorname{Fun}\operatorname{Op}_{{}^LG}(D^\times).
    \tag{8.3-1}
\end{equation}
本节讨论该同构的种种性质。

\subsection{$Z(\hat{\mathfrak g})$ 与 $\operatorname{Fun}\operatorname{Op}_{{}^LG}(D^\times)$ 之间的同构}

回想定理 8.2.5 建立的顶点 Poisson 代数的 $\operatorname{Aut}\mathcal O$-等变同构
\begin{equation}
    \mathfrak z(\hat{\mathfrak g})_{\kappa_0}\simeq\operatorname{Fun}\operatorname{Op}_{{}^LG}(D)_{\kappa^\vee_0}.
    \tag{8.3-2}
\end{equation}
对顶点 Poisson 代数 $P$，$U(P)=\operatorname{Lie}(P)$ 上的李代数结构给出交换代数 $\widetilde U(P)$ 上的 Poisson 代数结构。特别地，$\widetilde U(\mathfrak z(\hat{\mathfrak g})_{\kappa_0})=Z(\hat{\mathfrak g})$（见命题 4.3.4）是 Poisson 代数。

$Z(\hat{\mathfrak g})_{\kappa_0}$ 上相应的 Poisson 结构可描述如下。把 $\widetilde U_\kappa(\hat{\mathfrak g})$ 看作依赖参数 $\epsilon=(\kappa-\kappa_c)/\kappa_0$ 的结合代数单参数族 $A_\epsilon$。则 $Z(\hat{\mathfrak g})$ 是 $A_0$ 的中心。给定 $x,y\in Z(\hat{\mathfrak g})$，设 $\widetilde x,\widetilde y$ 是它们在 $A_\epsilon$ 中的提升。则 Poisson 括号 $\{x,y\}$ 定义为对易子 $[\widetilde x,\widetilde y]$ 视为 $\epsilon$ 的函数时的 $\epsilon$-线性项（它与提升的选取无关）。\footnote{这确实是 Poisson 结构，此事实最早由 V. Drinfeld 在 T. Hayashi 的工作 [1988] 之后注意到。}把带此 Poisson 结构的中心 $Z(\hat{\mathfrak g})$ 记作 $Z(\hat{\mathfrak g})_{\kappa_0}$。

同样，顶点 Poisson 代数 $\operatorname{Fun}\operatorname{Op}_{{}^LG}(D)_{\kappa^\vee_0}$ 给出拓扑 Poisson 代数 $\widetilde U(\operatorname{Fun}\operatorname{Op}_{{}^LG}(D)_{\kappa^\vee_0})$。按引理 4.3.5，后者以 $\operatorname{Aut}\mathcal O$-等变的方式同构于 $\operatorname{Fun}\operatorname{Op}_{{}^LG}(D^\times)$。因此同构 (8.3-2) 给出拓扑 Poisson 代数同构
\[
    Z(\hat{\mathfrak g})_{\kappa_0}\simeq\operatorname{Fun}\operatorname{Op}_{{}^LG}(D^\times)_{\kappa^\vee_0}.
\]
这里用下标 $\kappa^\vee_0$ 表示该 Poisson 结构对 ${}^L\mathfrak g$ 上内积 $\kappa^\vee_0$ 的依赖。下一节将给出该 Poisson 结构的另一定义。于是得到如下结果，它最初由 V. Drinfeld 猜想。\transnote{这里的「取包络」出现两次而含义不同：$\widetilde U(\mathfrak z(\hat{\mathfrak g})_{\kappa_0})=Z(\hat{\mathfrak g})$ 是把顶点 Poisson 代数取 $\operatorname{Lie}$ 再完备化，得到结合代数；$\widetilde U(\operatorname{Fun}\operatorname{Op}_{{}^LG}(D)_{\kappa_0^\vee})\simeq\operatorname{Fun}\operatorname{Op}_{{}^LG}(D^\times)$ 则是把「圆盘上的函数」完备化为「去心圆盘上的函数」（引理 4.3.5）。后者正是定理 8.3.1 中 $D$ 变成 $D^\times$ 的唯一原因。落点：§8.3.1，原书 p.246。[来源：本书 p.246，命题 4.3.4 与引理 4.3.5.]}

\begin{theorem}
    作为 Poisson 代数，中心 $Z(\hat{\mathfrak g})_{\kappa_0}$ 同构于 Poisson 代数 $\operatorname{Fun}\operatorname{Op}_{{}^LG}(D^\times)_{\kappa^\vee_0}$。此外，该同构是 $\operatorname{Aut}\mathcal O$-等变的。
\end{theorem}

% [Tutor Note: p.246 这里的两次“取包络”不可混淆：$\widetilde U(\mathfrak z(\hat{\mathfrak g})_{\kappa_0})=Z(\hat{\mathfrak g})$ 把顶点 Poisson 代数变成 Poisson 代数（取 $\operatorname{Lie}$ 再完备化），而 $\widetilde U(\operatorname{Fun}\operatorname{Op}_{{}^LG}(D)_{\kappa^\vee_0})\simeq\operatorname{Fun}\operatorname{Op}_{{}^LG}(D^\times)$ 说明“圆盘上的函数”完备化后成为“去心圆盘上的函数”。后者正是引理 4.3.5 的内容，也是定理 8.3.1 中 $D$ 变成 $D^\times$ 的唯一原因。]

\subsection{$\operatorname{Fun}\operatorname{Op}_G(D^\times)_{\upsilon_0}$ 上的 Poisson 结构}

上面我们从 $\operatorname{Fun}\operatorname{Op}_G(D)$ 上的顶点 Poisson 结构得到了代数 $\operatorname{Fun}\operatorname{Op}_G(D^\times)$ 上的 Poisson 结构，而后者又来自把 $\operatorname{Fun}\operatorname{Op}_G(D)$ 实现为 $\pi^\vee(\mathfrak g)_{\upsilon_0}$ 的顶点 Poisson 子代数（即经典 W-代数 $\mathbf W(\mathfrak g)_{\upsilon_0}$）。现在说明如何用称为 \emph{Drinfeld–Sokolov 约化}的 Hamilton 约化得到该 Poisson 结构 [Drinfeld and Sokolov 1985]。

从 $D^\times$ 上平凡 $G$-丛上联络的 Poisson 流形 $\operatorname{Conn}_{\mathfrak g}$ 出发，即形如 $\nabla=\partial_t+A(t)$（$A(t)\in\mathfrak g((t))$）的算子。该流形上的 Poisson 结构来自把它与仿射 Kac–Moody 代数 $\hat{\mathfrak g}_{\upsilon_0}$ 的对偶空间中的超平面等同，其中 $\upsilon_0$ 是 $\mathfrak g$ 上非零不变内积。确实，$\hat{\mathfrak g}_{\upsilon_0}$ 的拓扑对偶空间可等同于 $D^\times$ 上平凡丛上一切 $\lambda$-联络的空间，见 [Frenkel and Ben-Zvi 2004] §16.4。具体说，把 $\hat{\mathfrak g}_{\upsilon_0}=\mathfrak g((t))\oplus\mathbb C\mathbf 1$ 作向量空间分裂。则 $\lambda$-联络 $\partial_t+A(t)$ 给出 $\hat{\mathfrak g}_{\upsilon_0}$ 上的线性泛函，它在 $(B(t)+b\mathbf 1)\in\mathfrak g((t))\oplus\mathbb C\mathbf 1=\hat{\mathfrak g}_{\upsilon_0}$ 上取值
\[
    \lambda b+\operatorname{Res}\upsilon_0(A(t),B(t))dt.
\]
注意在此等同下，$\operatorname{Aut}\mathcal O$ 通过坐标变换的作用，与 $G((t))$ 在 $\lambda$-联络空间（相应地，$\hat{\mathfrak g}_{\upsilon_0}$ 的对偶空间）上的余伴随（相应地，规范）作用一致。空间 $\operatorname{Conn}_{\mathfrak g}$ 现在等同于 $\hat{\mathfrak g}^*_{\upsilon_0}$ 中由那些在中心元 $\mathbf 1$ 上取值 $1$ 的泛函构成的超平面。

对偶空间 $\hat{\mathfrak g}^*_{\upsilon_0}$ 带有称为 Kirillov–Kostant 结构的典范 Poisson 结构（见例如 [Frenkel and Ben-Zvi 2004] §16.4.1）。由于 $\mathbf 1$ 是中心元，该 Poisson 结构限制到我们已与 $\operatorname{Conn}_{\mathfrak g}$ 等同的超平面上。于是得到 $\operatorname{Conn}_{\mathfrak g}$ 上的 Poisson 结构。

群 $N((t))$ 通过规范变换作用在 $\operatorname{Conn}_{\mathfrak g}$ 上。该作用对应 $N((t))$ 在 $\hat{\mathfrak g}^*_{\upsilon_0}$ 上的余伴随作用，并且是 Hamilton 作用，矩映射为与嵌入 $\mathfrak n((t))\hookrightarrow\hat{\mathfrak g}_{\upsilon_0}$ 对偶的满射 $m:\operatorname{Conn}_{\mathfrak g}\to\mathfrak n((t))^*$。我们在 $\mathfrak n((t))^*$ 中取一个由线性泛函 $\psi$ 代表的单点余伴随 $N((t))$-轨道，$\psi$ 等于复合 $\mathfrak n((t))\to\mathfrak n/[\mathfrak n,\mathfrak n]((t))=\bigoplus_{i=1}^{\ell}\mathbb C((t))\cdot e_i$ 与泛函
\[
    (x_i(t))_{i=1}^{\ell}\mapsto\sum_{i=1}^{\ell}\operatorname{Res}_{t=0}x_i(t)dt
\]
的复合。与引理 4.2.2 的证明同样可证，$N((t))$ 在 $m^{-1}(\psi)$ 上的作用自由。此外，商 $m^{-1}(\psi)/N((t))$（即 Poisson 约化流形）典范等同于 $D^\times$ 上 $G$-opers 的空间。于是得到拓扑函数代数 $\operatorname{Fun}\operatorname{Op}_G(D^\times)$ 上的 Poisson 结构。

这样，对 $\mathfrak g$ 上每个非零不变内积 $\upsilon_0$，我们现在有 $\operatorname{Fun}\operatorname{Op}_G(D^\times)$ 上的两个 Poisson 结构。\transnote{两个 Poisson 结构来源不同：一来自 §8.1.4 的水平形变极限（顶点 Poisson 代数取包络），二来自 Drinfeld–Sokolov 约化（把 Kirillov–Kostant 结构限制到 $\mathbf 1=1$ 的超平面再约化）。引理 8.3.2 通过 $C^\bullet_\infty(\mathfrak g)$ 的零次上同调把两者都认成 $\mathbf W(\mathfrak g)_{\upsilon_0}$ 的包络，故二者相等。落点：§8.3.2，原书 p.247–248。[来源：本书 p.247–248；Frenkel–Ben-Zvi 2004, §15.4 与 §16.8；Drinfeld–Sokolov, \emph{J. Sov. Math.} 30 (1985).]}

\begin{lemma}
    这两个 Poisson 结构相同。
\end{lemma}

\begin{proof}
在 [Frenkel and Ben-Zvi 2004] §15.4 与 §16.8 中，我们定义了复形 $C^\bullet_\infty(\mathfrak g)$，并证明其零次上同调典范同构于 $\mathbf W(\mathfrak g)_{\upsilon_0}$（带 §8.1.4 引入的顶点 Poisson 代数结构），且其余上同调为零。该复形是计算上述 Drinfeld–Sokolov 约化结果的 BRST 复形的顶点 Poisson 代数版本。特别地，我们把带由 Drinfeld–Sokolov 约化所得 Poisson 结构的 $\operatorname{Fun}\operatorname{Op}_G(D^\times)$ 与
\[
    \widetilde U(\mathbf W(\mathfrak g)_{\upsilon_0})
\]
等同，其中 Poisson 结构对应上述 $\mathbf W(\mathfrak g)_{\upsilon_0}$ 上的顶点 Poisson 结构。
\end{proof}

% [Tutor Note: p.247 引理 8.3.2 的证明把两个 Poisson 结构“都认成”同一个顶点 Poisson 代数 $\mathbf W(\mathfrak g)_{\upsilon_0}$ 的包络：一个来自 §8.1.4 的水平形变极限，一个来自 Drinfeld–Sokolov 的 Hamilton 约化。所需的技术输入是复形 $C^\bullet_\infty(\mathfrak g)$ 的零次上同调计算（[Frenkel and Ben-Zvi 2004] 定理 15.4.12 一类的结论），此处只以一句话带过，读者若追细节须回到该书第 15、16 章。]

\subsection{Miura 变换作为 Harish--Chandra 同态}

我们在 §5.1 已讨论过的 Harish--Chandra 同态，是从 $U(\mathfrak g)$ 的中心 $Z(\mathfrak g)$（$\mathfrak g$ 为单李代数）到 $\mathfrak h^*$ 上多项式代数 $\operatorname{Fun}\mathfrak h^*$ 的同态。它把 $Z(\mathfrak g)$ 与 $\mathfrak h^*$ 上 $W$-不变多项式代数 $(\operatorname{Fun}\mathfrak h^*)^W$ 等同。为构造此同态，需给每个 $\lambda\in\mathfrak h^*$ 配一个中心特征标。该中心特征标就是中心在 Verma 模 $M_{\lambda-\rho}$ 上的作用特征标。

仿射情形，我们构造类似的同态，从 $\widetilde U_{\kappa_c}(\hat{\mathfrak g})$ 的中心 $Z(\hat{\mathfrak g})$ 到 $D^\times$ 上 ${}^LH$-丛 $\Omega^{-\rho}$ 上联络空间上的函数代数 $\operatorname{Fun}\operatorname{Conn}(\Omega^{-\rho})_{D^\times}$。按推论 6.1.5，$\operatorname{Conn}(\Omega^{-\rho})_{D^\times}$ 的点参数化临界水平的 Wakimoto 模。于是对每个 $\overline\nabla\in\operatorname{Conn}(\Omega^{-\rho})_{D^\times}$ 有临界水平的 Wakimoto 模 $W_{\overline\nabla}$。下面的定理描述 Harish--Chandra 同态的仿射类比。

注意 $\operatorname{Fun}\operatorname{Conn}(\Omega^{-\rho})_{D^\times}$ 是多项式代数 $\mathbb C[b_{i,n}]$（$i=1,\dots,\ell$；$n\in\mathbb Z$）关于下述拓扑的完备化：$0$ 的开邻域基由 $b_{i,n}$（$n<N$）生成的理想构成。与引理 4.3.5 的证明同样可证，它同构于
\[
    \widetilde U(\operatorname{Fun}\operatorname{Conn}(\Omega^{-\rho})_D).
\]
我们还把拓扑代数 $\operatorname{Fun}\operatorname{Conn}(\Omega^{\rho})_{D^\times}$ 定义为 $\widetilde U(\operatorname{Fun}\operatorname{Conn}(\Omega^{\rho})_D)$。它是多项式代数 $\mathbb C[u_{i,n}]$（$i=1,\dots,\ell$；$n\in\mathbb Z$）关于下述拓扑的完备化：$0$ 的开邻域基由 $u_{i,n}$（$n<N$）生成的理想构成。同构 (8.2-14) 给出拓扑代数同构
\[
    \operatorname{Fun}\operatorname{Conn}(\Omega^{-\rho})_{D^\times}\longrightarrow\operatorname{Fun}\operatorname{Conn}(\Omega^{\rho})_{D^\times},
\]
在其下 $b_{i,n}\mapsto-u_{i,n}$。

自然遗忘态射 $\operatorname{MOp}_{{}^LG}(D^\times)_{\mathrm{gen}}\to\operatorname{Op}_{{}^LG}(D^\times)$，连同按命题 8.2.2 同样方式构造的等同
\[
    \operatorname{MOp}_{{}^LG}(D^\times)_{\mathrm{gen}}\simeq\operatorname{Conn}(\Omega^{\rho})_{D^\times},
\]
给出映射
\begin{equation}
    \mu:\operatorname{Conn}(\Omega^{\rho})_{D^\times}\longrightarrow\operatorname{Op}_{{}^LG}(D^\times).
    \tag{8.3-3}
\end{equation}
这是去心圆盘 $D^\times$ 上的 Miura 变换，类似于圆盘 $D$ 上的 Miura 变换 (8.2-3)。

回想推论 6.3.2 中描述的半无限抛物诱导的一般构造。把它用于 Borel 子代数 $\mathfrak b\subset\mathfrak g$ 的情形。设 $M$ 是顶点代数 $\mathbb M_{\mathfrak g}$ 上的模，等价地，是 Weyl 代数 $\mathcal A^{\mathfrak g}$ 上的模。设 $R$ 是交换顶点代数 $\pi_0=\operatorname{Fun}\operatorname{Conn}(\Omega^{-\rho})_D$ 上的模，等价地，是交换代数 $\operatorname{Fun}\operatorname{Conn}(\Omega^{-\rho})_{D^\times}$ 上的光滑模。则定理 6.1.6 的同态
\[
    w_{\kappa_c}:V_{\kappa_c}(\mathfrak g)\longrightarrow\mathbb M_{\mathfrak g}\otimes\pi_0
\]
在张量积 $\mathbb M_{\mathfrak g}\otimes R$ 上给出 $V_{\kappa_c}(\mathfrak g)$-模结构，等价地，给出光滑 $\hat{\mathfrak g}_{\kappa_c}$-模结构。

现在由定理 8.2.5 得到如下结果。

\begin{theorem}
    $Z(\hat{\mathfrak g})$ 在 $\hat{\mathfrak g}_{\kappa_c}$-模 $M\otimes R$ 上的作用不依赖 $M$，并经由同态 $Z(\hat{\mathfrak g})\to\operatorname{Fun}\operatorname{Conn}(\Omega^{-\rho})_{D^\times}$ 与 $\operatorname{Fun}\operatorname{Conn}(\Omega^{-\rho})_{D^\times}$ 在 $R$ 上的作用分解。相应的态射 $\operatorname{Conn}(\Omega^{-\rho})_{D^\times}\to\operatorname{Spec}Z(\hat{\mathfrak g})$ 嵌入交换图
    \begin{equation}
    \begin{CD}
    \operatorname{Conn}(\Omega^{-\rho})_{D^\times} @>{\sim}>> \operatorname{Conn}(\Omega^{\rho})_{D^\times}\\
    @VVV @VVV\\
    \operatorname{Spec}Z(\hat{\mathfrak g}) @>{\sim}>> \operatorname{Op}_{{}^LG}(D^\times)
    \end{CD}
    \tag{8.3-4}
    \end{equation}
    其中右竖箭头是 $D^\times$ 上的 Miura 变换 (8.3-3)。

    特别地，若 $R=\mathbb C_{\overline\nabla}$ 是对应于点 $\overline\nabla\in\operatorname{Conn}(\Omega^{-\rho})_{D^\times}$ 的一维模，则中心 $Z(\hat{\mathfrak g})$ 通过对应 $\overline\nabla$ 的 Miura 变换的中心特征标，作用在相应的 Wakimoto 模 $W_{\overline\nabla}=\mathbb M_{\mathfrak g}\otimes\mathbb C_{\overline\nabla}$ 上。
\end{theorem}

\begin{proof}
$U(V_{\kappa_c}(\mathfrak g))$ 从而 $\widetilde U_{\kappa_c}(\hat{\mathfrak g})$ 在 $W_{\overline\nabla}$ 上的作用，通过顶点代数同态 $V_{\kappa_c}(\mathfrak g)\to\mathbb M_{\mathfrak g}\otimes\pi_0(\mathfrak g)$ 得到。特别地，$Z(\hat{\mathfrak g})$ 在 $W_{\overline\nabla}$ 上的作用通过交换顶点代数同态 $\mathfrak z(\hat{\mathfrak g})\to\pi_0(\mathfrak g)$ 得到。但按 §6.2.5，$\pi_0(\mathfrak g)=\operatorname{Fun}\operatorname{Conn}(\Omega^{-\rho})_D$。因此定理的陈述可由定理 8.2.5 施加 §3.2.3 引入的包络代数函子 $V\mapsto\widetilde U(V)$ 得到。
\end{proof}

于是我们看到，Harish--Chandra 同态的仿射类比，是图 (8.3-4) 的右竖箭头；它不是别的，正是 Langlands 对偶群的 Miura 变换！特别地，它的像（在有限维情形由 $\mathfrak h^*$ 上 $W$-不变多项式构成）可描述为 W-代数屏蔽算子核之交。\transnote{这是经典 Harish--Chandra 同构的仿射类比：$Z(\mathfrak g)\to(\operatorname{Fun}\mathfrak h^*)^W$ 的像由 $W$-不变多项式刻画，而 (8.3-4) 的右竖箭头即对偶群的 Miura 变换，其像由「W-代数屏蔽算子核之交」刻画。两者是同一件事在 $D^\times$ 与 $D$ 上的两种表述：Wakimoto 模 $W_{\overline\nabla}$ 的中心特征标正是 $\overline\nabla$ 的 Miura 变换。落点：§8.3.3，(8.3-3)(8.3-4)，原书 p.249。[来源：本书 p.249；推论 6.1.5 与定理 6.1.6.]}

% [Tutor Note: p.249 全章收束于这一句：经典 Harish--Chandra 同态 $Z(\mathfrak g)\to(\operatorname{Fun}\mathfrak h^*)^W$ 的像由 $W$-不变多项式刻画，而仿射类比 $Z(\hat{\mathfrak g})\to\operatorname{Fun}\operatorname{Conn}(\Omega^{-\rho})_{D^\times}$ 的像由“W-代数屏蔽算子的核之交”刻画——两者是同一件事在 $D$ 与 $D^\times$ 上的两种表述。这也正是 Drinfeld 猜想“中心 = opers 上函数”的最终形式：Miura 变换把 Wakimoto 模的中心特征标送到 oper 参数。]
